% Research contribution of agent /root/monodromy.  Input into a parent document.
% Status is local to the assertions below; this is not a claim to solve all of A.3--A.5.
\section{A.3--A.5: proved statements and remaining gates}
\label{sec:monodromy-contribution}
\status{Incomplete universal determination; the individually labelled results below are proved.  Lead A's full-linear assertion is refuted, its positivity assertion is proved with an explicit threshold, and its plane-curve point-count consequence is proved.}

\subsection{Conventions and the generic rank for projective hypersurfaces}

Let $S_d=\mathbf P(H^0(\mathbf P^r,\mathcal O(d)))$, and let
$\pi:\mathcal Y_d\to S_d$ be the universal degree-$d$ hypersurface.
For $r\geq2$, the sheaf $\mathcal H=R^{r-1}\pi_*\mathcal O_{\mathcal Y_d}$
is locally free of rank
\[
 h_d=\begin{cases}\binom{d-1}{r}&d\geq r+1,\\0&d\leq r.\end{cases}
\]
This follows from the universal divisor exact sequence and the cohomology
of line bundles on projective space; formation of $\mathcal H$ commutes
with all base changes.  Absolute Frobenius induces a map
$F_S^*\mathcal H\to\mathcal H$.  On its invertibility locus $W_d$, the
kernel of $F-1$ is a locally constant $\mathbf F_p$-sheaf of rank $h_d$.
On the smooth locus it is $R^{r-1}\pi_*\mathbf F_p$, by the
Artin--Schreier sequence.  For $X=\mathbf P^r$, this is the vanishing
local system in the question, since $H^{r-1}(\mathbf P^r,\mathcal O)=0$.

\begin{proposition}[All degrees, all characteristics]\label{prop:projective-generic-hw-rank}
The generic degree-$d$ hypersurface of $\mathbf P^r$ has bijective
Frobenius on $H^{r-1}(\mathcal O)$, and its generic stable rank is $h_d$.
In particular, for plane curves the generic rank is
$g_d=(d-1)(d-2)/2$, and the smallest positivity threshold is $d_0=3$.
\end{proposition}

\begin{proof}
Work over an algebraic closure and take $d$ hyperplanes in general
position.  Their union $D$ is a simple normal crossings divisor.
The alternating resolution of $\mathcal O_D$ by the structure sheaves
of intersections has only $H^0$ in its first hypercohomology page,
because every nonempty intersection is a projective space.  Frobenius
is bijective on each of these $H^0$'s and hence on the cohomology of
their complex.  It is therefore bijective on
$H^{r-1}(D,\mathcal O_D)$.  Thus $W_d$ is nonempty.  The smooth locus
is a nonempty dense open of the irreducible scheme $S_d$, so it meets
$W_d$.  Invertibility is open, which proves the assertion.  This proof
establishes only bijectivity on the specified coherent cohomology;
it does not conflate this condition with every formulation of
ordinarity for higher dimensional varieties.
\end{proof}

\par\smallskip\noindent\textbf{Comparison with Illusie.}
Illusie (1990), Theorem~0.1, p.~375, proves the stronger statement that
the general smooth complete intersection of every specified dimension
and multidegree is ordinary: its Newton and Hodge polygons coincide.
The theorem and the explanation of the earlier Koblitz argument on
pp.~375--376 were inspected in the publisher's two-page preview of
the reprint.  In particular, generic Hasse--Witt bijectivity follows;
the proposition above gives a separate complete proof of exactly the
part needed here.  The parameter space in Illusie's theorem varies all
defining equations.  It does not, by restriction of an open set alone,
prove generic ordinarity for hypersurface sections of every fixed $X$.

\subsection{The determinant character and the Hasse invariant}
\label{subsec:hasse-determinant}

Let $S$ be connected and let $(\mathcal E,F)$ be a vector bundle of
rank $h$ in characteristic $p$, with $F:F_S^*\mathcal E\simeq\mathcal E$.
Write $V=\ker(F-1)$ and $\mathcal L=\det\mathcal E$.  Taking
determinants gives an invertible section
\[
 \operatorname{Ha}_F\in
 H^0\bigl(S,\operatorname{Hom}(F_S^*\mathcal L,\mathcal L)\bigr)
 =H^0(S,\mathcal L^{\otimes(1-p)}).
\]
This is the Hasse invariant of the specified Frobenius module.  It
determines the determinant character exactly, including its arithmetic
part: $\det V$ is the rank-one Lang local system
$\ker(F_{\mathcal L}-1)$.  If $e$ trivializes $\mathcal L$ and
$F_{\mathcal L}(e)=h e$, a fixed nonzero vector $ze$ satisfies
\begin{equation}\label{eq:determinant-kummer}
 z^{p-1}=h^{-1}.
\end{equation}
Thus the $\mathbf F_p^\times$-torsor of frames of $\det V$ is the
Kummer torsor \eqref{eq:determinant-kummer}; a deck transformation
$z\mapsto az$ acts on the determinant fibre by $a$.  Changing the
trivialization of $\mathcal L$ changes this equation by a
$(p-1)$st power, as required.  Over $\mathbf F_q$, geometric Frobenius
acts on $\det V$ by the determinant of the $q$-power coherent
Frobenius; arithmetic Frobenius is its inverse.  This description
holds for the stable bundle of the relative vanishing system on its
constant-rank locus as well as for the full ordinary coherent module.
It does not assert that the image of the Kummer character is always
all of $\mathbf F_p^\times$: that depends on its Kummer class.

\subsection{The exact plane-curve monodromy in degrees at most four}
\label{subsec:small-plane-monodromy}

For $d=1,2$, the rank is zero and both monodromy groups are trivial.
For $d=3$, both images are $\mathrm{GL}_1(\mathbf F_p)$: for odd $p$
this follows from the Legendre subfamily and its simple Hasse roots,
proved below; for $p=2$ the assertion means the trivial group on a
one-dimensional space.  A nonzero rank-one representation is
irreducible.

For $d=4$, both images are $\mathrm{GL}_3(\mathbf F_p)$, and the
representation is irreducible.  Here is why the moduli-space theorem
actually applies.  Every nonhyperelliptic smooth genus-three curve is
canonically a smooth plane quartic, and conversely.  The stack quotient
of the parameter space of smooth plane quartics by $\mathrm{PGL}_3$
is the nonhyperelliptic open substack of $\mathcal M_3$.
The map from the parameter space to this quotient is a
$\mathrm{PGL}_3$-torsor in the stack sense, with geometrically connected
fibres, and is surjective on etale fundamental groups: a connected
finite etale cover of the quotient remains connected after pullback
along a torsor under a connected algebraic group.  Restriction to
the ordinary locus preserves this assertion.  That locus is a
nonempty open of the ordinary locus of $\mathcal M_3$, hence its
fundamental group has the same finite monodromy image.
Achter--Pries (2008), Proposition~4.8, says that this image on the
etale $p$-torsion of the Jacobian is $\mathrm{GL}_3(\mathbf F_p)$.
Dualizing gives the same full group on $V$.
The arithmetic image contains the geometric image and is already
a subgroup of the same finite general linear group, so it too is full.
No analogous moduli-space surjectivity assertion is used for degrees
$d\geq5$.  In characteristic two the separate contact-curve argument
and group induction instead determine every degree, in
Theorem~\ref{thm:plane-char2-all-groups}.

\subsection{A universal obstruction in characteristic two}
\label{subsec:odd-plane-obstruction}

Let $p=2$ and let $d\geq3$ be odd.  Write
\[
 f=\sum_{|a|=d}c_aX^{a_1}Y^{a_2}Z^{a_3},\qquad
 I_d=\{u\in\mathbf Z_{>0}^3:|u|=d\}.
\]
Put $c_a=0$ if a coordinate of $a$ is negative.  On a coefficient
chart of $S_d$, use the Cech basis of
$H^2(\mathbf P^2,\mathcal O(-d))$ indexed by $I_d$.
The matrix of the semilinear Frobenius on $H^1(C_f,\mathcal O)$ is
\begin{equation}\label{eq:plane-hw-char2}
 H_{u,v}=c_{2v-u}\quad(u,v\in I_d),\qquad F(x)=Hx^{[2]}.
\end{equation}
Indeed, the connecting map from
$0\to\mathcal O(-d)\xrightarrow{f}\mathcal O\to\mathcal O_{C_f}\to0$
identifies Frobenius with multiplication by $f$ after squaring the
Cech representative $X^{-v_1}Y^{-v_2}Z^{-v_3}$.

Define a column $w=(w_u)$ by
\[
 w_u=\begin{cases}c_u&u_1,u_2,u_3\text{ all odd},\\0&\text{otherwise}.
 \end{cases}
\]
Then the following is an identity over the polynomial ring in the
coefficients, not merely a numerical identity at selected curves:
\begin{equation}\label{eq:universal-covector}
 H^t w=w^{[2]}.
\end{equation}
For its $v$-coordinate, pair the term $c_{2v-u}c_u$ indexed by $u$
with the term indexed by $2v-u$.  A nonzero term has both indices
positive and all odd, so the involution preserves the indexing set.
All two-element orbits cancel in characteristic two.  The only possible
fixed index is $u=v$, contributing $c_v^2$ exactly when $v$ is all odd.

If $H$ is invertible, then $w\ne0$.  Otherwise all coefficients with
three odd exponents vanish, and every row with all-odd index $u$
in \eqref{eq:plane-hw-char2} vanishes.  Such an index exists because
$(1,1,d-2)\in I_d$.  This contradicts invertibility.

For a Lang fixed vector $x$, namely $Hx^{[2]}=x$, equation
\eqref{eq:universal-covector} gives
\[
 (w^tx)^2=(w^{[2]})^tx^{[2]}=w^tHx^{[2]}=w^tx.
\]
Thus $x\mapsto w^tx$ is an $\mathbf F_2$-valued linear functional on
the local system $V=\ker(F-1)$.  It is nonzero, since the Lang fixed
vectors form a basis after an etale cover.  It is a section of
$V^\vee$, so it is fixed by both geometric and arithmetic monodromy.
Working on a dense coefficient chart suffices: its intersection with
the ordinary smooth locus is dense, and the map on fundamental groups
from a nonempty open of a normal integral scheme is surjective.

\begin{theorem}[Refutation of full linear monodromy]\label{thm:odd-plane-char2-obstruction}
For every finite field of characteristic two and every odd $d\geq5$,
let $\bar\eta$ be a geometric generic point of the ordinary smooth
degree-$d$ plane-curve locus $U_d$, and put
$V=H^1_{\mathrm{et}}(C_{\bar\eta},\mathbf F_2)$.
Both images
\[
 G_{\mathrm{geom}}=\mathrm{im}\bigl(\pi_1(U_d\otimes\bar k,\bar\eta)
 \to\mathrm{GL}(V)\bigr),\qquad
 G_{\mathrm{arith}}=\mathrm{im}\bigl(\pi_1(U_d,\bar\eta)
 \to\mathrm{GL}(V)\bigr)
\]
fix a nonzero covector.  Consequently both are contained in the
proper subgroup
\[
 \mathrm{Stab}(\ell)\simeq
 \mathbf F_2^{\,g_d-1}\rtimes\mathrm{GL}_{g_d-1}(\mathbf F_2),
\]
and $V$ is reducible, with invariant hyperplane $\ker\ell$.
There is no integer $d_1(\mathbf P^2)$ for which the proposed full
linear geometric or arithmetic monodromy holds for every $d\geq d_1$
in characteristic two.
\end{theorem}

The theorem proves containment, not equality with the displayed
stabilizer.  Determining the exact image remains a separate gate.
The determinant character in this characteristic is automatically
trivial, because $\mathbf F_2^\times=\{1\}$.

\par\smallskip\noindent\textbf{Geometric interpretation and literature check.}
The obstruction agrees with the difference between the effective
geometric theta characteristic
$\mathcal O_C((d-3)/2)$ and the canonical characteristic-two theta
characteristic, which is non-effective for an ordinary curve.
Cais, Ellenberg, and Zureick-Brown (2013, Theorem~4.2) prove the resulting
rational two-torsion statement for odd degree plane curves.
Castryck, Streng, and Testa (2014, Corollary~13, Remark~14, and
Theorem~16) give the ordinary and Hasse--Witt formulations.
The explicit identity above avoids any descent ambiguity for theta
characteristics over imperfect coefficient fields.

\par\smallskip\noindent\textbf{Explicit checked fibre.}
Over $\mathbf F_2$, take
\begin{align*}
f={}&x^5+x^4y+x^3y^2+x^2y^3+x^2y^2z+xy^4+xyz^3+xz^4\\
 &+y^4z+y^3z^2+yz^4.
\end{align*}
In the ordered basis
$(1,1,3),(1,2,2),(1,3,1),(2,1,2),(2,2,1),(3,1,1)$, its matrix is
\[
H=\begin{pmatrix}
1&0&0&0&0&0\\1&0&1&0&1&1\\0&1&0&0&0&0\\
1&1&0&0&1&1\\0&0&1&0&1&0\\0&0&0&1&0&0
\end{pmatrix},\qquad w=(1,0,0,0,0,0)^t.
\]
Its rank is six.  The ideal $(f,f_x,f_y,f_z)$ has Groebner basis
$\{1\}$ on each of the three standard projective charts, proving
geometric smoothness.  Direct enumeration over
$\mathbf F_{2^m}$ for $m=1,\ldots,6$ gives respectively
$3,3,9,23,38,33$ points; the traces of $H^m$ modulo two are
$0,0,0,0,1,0$.  The program \texttt{code/plane\_char2.py} and its JSON
output record the input, field arithmetic, certificates, and timings.
The fibre is a check of the formulas; the universal identity, rather
than this one fibre, proves the generic monodromy obstruction.

\subsection{Plane curves: every point-count residue in every open}
\label{subsec:plane-residues}

\begin{theorem}\label{thm:plane-all-residues}
Let $k$ be a finite field of characteristic $p$.  For every $d\geq4$,
every nonempty Zariski open $U\subset S_d$, and every $a\in\mathbf F_p$,
there exist a finite extension $k'/k$ and a smooth plane curve
$C\in U(k')$ of degree $d$ with $\#C(k')\equiv a\pmod p$.
For $d=1,2$ the only residue for smooth curves is $1$.
For $d=3$, the residues on the ordinary smooth locus are exactly
$\mathbf F_p\setminus\{1\}$, and every nonempty open of that locus
attains every such residue after a finite extension.  Supersingular
smooth cubics give residue $1$, so all residues occur among all smooth
cubics; the every-open assertion for all residues fails in degree three.
\end{theorem}
\begin{proof}

\par\smallskip\noindent\textbf{Elliptic input.}
For every $\lambda\in\mathbf F_p^\times$, there is an ordinary elliptic
curve $E/\mathbf F_Q$, with $\mathbf F_Q/k$ finite, for which geometric
Frobenius on $H^1_{\mathrm{et}}(E_{\bar k},\mathbf F_p)$ is multiplication
by $\lambda$.  For $p=2$ take any ordinary elliptic curve.  For odd $p$,
write $m=(p-1)/2$.  The Legendre family
$y^2=x(x-1)(x-t)$ has Hasse polynomial
\[
 A(t)=(-1)^m\sum_{i=0}^{m}\binom mi^2t^i.
\]
Its roots are simple and avoid $0,1$: its values there are nonzero,
and it satisfies
$t(1-t)A''+(1-2t)A'-m^2A=0$ modulo $p$.
A repeated root away from $0,1$ would, by successive differentiation,
force all derivatives through degree $m<p$ to vanish, a contradiction.
At a simple root, the Kummer cover $z^{p-1}=A(t)^{-1}$ is totally
ramified of degree $p-1$.  Its geometric monodromy is therefore the
whole $\mathbf F_p^\times$.  Finite-field Chebotarev on the ordinary
Legendre locus gives every scalar $\lambda$.
Moreover we may replace $Q$ by $Q^e$ with arbitrarily large
$e\equiv1\pmod{p-1}$ without changing $\lambda$.

\par\smallskip\noindent\textbf{A nodal curve with prescribed residue.}
Put $r=d-3\geq1$ and embed $E$ as a plane cubic by $|3O|$.
Choose rational lines $L_1,\ldots,L_r$ in general position such that
all intersections are transverse, no three components meet, the
intersection $E\cap L_i$ is split for $i\geq2$, and
$E\cap L_1$ has either cycle type $3$ or cycle type $1+2$ under
Frobenius.  Write $t=0$ or $1$ for the number of rational points in
this last intersection.  Such lines exist after an extension of the
above kind, as follows.

Split lines arise from distinct $P,Q,R\in E(\mathbf F_Q)$ with
$P+Q+R=0$, giving order $Q^2$ choices.  Lines of type $1+2$ arise from
$P\in E(\mathbf F_{Q^2})\setminus E(\mathbf F_Q)$ and
$R=-(P+\pi P)\in E(\mathbf F_Q)$; the three points are distinct,
collinear, and Frobenius invariant.  There are order $Q^2$ choices.
For type $3$, use nonrational points in
$\ker(1+\pi+\pi^2)$.  This is the kernel of a separable isogeny,
is contained in $E(\mathbf F_{Q^3})$, and has cardinality
$\#E(\mathbf F_{Q^3})/\#E(\mathbf F_Q)$, of order $Q^2$.
At most nine of these points are rational, since rational points in
the kernel lie in $E[3]$.  Thus there are order $Q^2$ suitable lines
of each required type.  The finitely many incidence and tangency
conditions excluded at each choice remove only order $Q$ rational
lines.  This proves existence for sufficiently large $Q$.

The connected nodal curve
$C_0=E\cup L_1\cup\cdots\cup L_r$ has degree $d$.
Its normalization sequence
\[
0\longrightarrow\mathcal O_{C_0}\longrightarrow
\mathcal O_E\oplus\bigoplus_i\mathcal O_{L_i}
\longrightarrow\bigoplus_{P\text{ node}}k(P)\longrightarrow0
\]
shows that Frobenius on $H^1(C_0,\mathcal O_{C_0})$ is bijective:
it is bijective on the finite reduced node scheme, on constants of
the normalization, and on $H^1(E,\mathcal O_E)$.
Inclusion--exclusion gives
\begin{align*}
 \#C_0(\mathbf F_Q)
 &\equiv (1-\lambda)+r-\{t+3(r-1)\}-\binom r2\\
 &\equiv 4-2r-\binom r2-t-\lambda\pmod p.
\end{align*}
As $t\in\{0,1\}$ and $\lambda\in\mathbf F_p^\times$ vary, these
residues fill $\mathbf F_p$, including when $p=2$.

\par\smallskip\noindent\textbf{Passage to an arbitrary smooth open.}
The ordinary coherent-cohomology locus $W_d\subset S_d$ introduced
above is geometrically integral and contains $C_0$.
The finite etale local system $\ker(F-1)$ is defined over all of $W_d$,
including its singular fibres.  After base changing to $\mathbf F_Q$,
the Frobenius class of $C_0$ belongs to the finite arithmetic
monodromy group.  Chebotarev gives that same class at a closed point
of any nonempty open of $W_d$.  Apply this to
$U\cap W_d\cap S_d^{\mathrm{smooth}}$.
The coherent Frobenius trace formula
\[
 \#C(\mathbf F_{Q^e})\equiv
 1-\operatorname{Tr}\bigl(\operatorname{Frob}_{Q^e}^{\mathrm{geom}},
 H^1_{\mathrm{et}}(C_{\bar k},\mathbf F_p)\bigr)\pmod p
\]
also applies to the proper nodal curve $C_0$, so equality of the
classes gives the desired residue.  More precisely, the class at
$C_0$ maps to degree one in the finite constant-field quotient, and
Chebotarev supplies it at degrees congruent to one in that quotient;
the permitted residue fields are finite extensions of $k$.
This argument never assumes full general linear monodromy for
degree-$d$ plane curves.
\end{proof}

\par\smallskip\noindent\textbf{Primary-source verification of the two inputs.}
Katz--Sarnak (1999), Theorem~9.7.13, p.~290, is the finite-monodromy
equidistribution theorem, with the constant-field coset condition made
explicit in 9.7.1--9.7.9.  To apply its characteristic-zero formulation
to our finite $\mathbf F_p$-representation, take a faithful finite
dimensional characteristic-zero representation of its finite image
(for example the regular representation), choose $\ell\ne p$, and
regard it as a lisse $\overline{\mathbf Q}_\ell$-sheaf.  It is pure of
weight zero, and its finite geometric and arithmetic groups have the
required normality.  The theorem gives a positive limiting proportion
for every nonempty conjugacy class in the required coset.  Restriction
to a nonempty open preserves the finite monodromy image because the
base is normal.  Katz (1973), SGA~7~II, Expose~XXII, Corollary~3.2,
formulas (3.2.1)--(3.2.2), pp.~410--411, gives precisely the trace
congruence for every proper scheme over a finite field, including
our nodal curve.  Proposition~2.0 and formula (2.2.4) of the same
expose identify its coherent and etale formulations.

\subsection{An explicit positivity threshold for every smooth ambient variety}
\label{subsec:explicit-positivity}

Here $X\subset\mathbf P^N$ is smooth projective of pure dimension
$n\geq2$ over a perfect field of characteristic $p$, and $L=\mathcal O_X(1)$.
Geometric connectedness is not needed for this proposition.  All
dimensions in the following definitions may be computed after extending
the ground field to an algebraic closure.  Put
\[
 D=\deg_L X,\qquad
 a=\dim H^{n-1}(X,\mathcal O_X)_{\mathrm{ss}},
\]
and, for $1\leq m\leq n-2$, define the explicit cohomological bounds
\[
 B_m=\sum_{j=0}^{m}\binom mj
       h^{n-m-1+j}(X,L^{-j}),\qquad
 B=D+a+\sum_{m=1}^{n-2}B_m.
\]
Let $r_\omega$ be any Castelnuovo--Mumford regularity bound for
$i_*\omega_X$ on $\mathbf P^N$, and put
\begin{equation}\label{eq:explicit-d0}
 d_0=\max\left\{2n,1,r_\omega-1,
       \left\lfloor\frac{2^n n!B}{D}\right\rfloor+1\right\}.
\end{equation}
Every term is determined by the embedding and finitely many coherent
cohomology groups of $X$; this is a computable bound, not an unspecified
``sufficiently large'' constant.

\begin{theorem}\label{thm:all-X-explicit-positive-rank}
For every $d\geq d_0$, the generic smooth hypersurface section $Y$ of
degree $d$ has
$\dim E(Y,\mathcal O_Y)_{\mathrm{ss}}>0$.
More precisely, for every $d\geq\max\{1,r_\omega-1\}$, its generic
rank $f_{X,d}$ satisfies
\begin{equation}\label{eq:generic-rank-lower-bound}
 f_{X,d}\ \geq\
 D\binom dn-D\binom d{n-1}
 -\sum_{m=1}^{n-2}B_m\binom dm-a,
\end{equation}
where a binomial coefficient with lower index greater than $d$ is zero.
Consequently
\[
 f_{X,d}=\frac{D}{n!}d^n+O_X(d^{n-1})\quad(d\longrightarrow\infty).
\]
\end{theorem}

\begin{proof}
Take $d$ hyperplanes in simultaneous general position over the
algebraic closure and set $Y_0=\bigcup_{i=1}^d(X\cap H_i)$.
Every intersection $Z_I=X\cap\bigcap_{i\in I}H_i$ with $|I|\leq n$
is smooth of dimension $n-|I|$, and intersections of $n+1$ components
are empty.  Intersections with $|I|=n$ consist of $D$ reduced points.
The alternating normal-crossings resolution gives the Frobenius
equivariant spectral sequence
\[
 E_1^{s,t}=\bigoplus_{|I|=s+1}H^t(Z_I,\mathcal O_{Z_I})
 \ \Longrightarrow\ H^{s+t}(Y_0,\mathcal O_{Y_0}).
\]
The term $E_1^{n-1,0}$ has dimension $D\binom dn$ and bijective
Frobenius.  It has no outgoing differential on any page.  The image
of the incoming $d_1$ has dimension at most
$D\binom d{n-1}$: each smooth positive dimensional intersection has
at most $D$ connected components, since each has positive integral
degree.  For $r=2,\ldots,n-1$, the possible incoming $d_r$ has source
$E_r^{n-1-r,r-1}$.  Writing $m=n-r$, its dimension is at most
$\binom dm B_m$.  Indeed, the Koszul resolution of a transverse
$m$-fold hyperplane section bounds
\[
 h^{n-m-1}(Z_I,\mathcal O_{Z_I})
 \leq\sum_{j=0}^{m}\binom mj
 h^{n-m-1+j}(X,L^{-j})=B_m.
\]
Every quotient of the top term by an incoming image still has
bijective Frobenius: a Frobenius-stable subspace of a finite dimensional
bijective semilinear space over a perfect field itself has bijective
Frobenius.  Hence $E_\infty^{n-1,0}$ is a bijective-Frobenius subquotient
of $H^{n-1}(Y_0,\mathcal O_{Y_0})$ of dimension at least the first three
terms of the right side of \eqref{eq:generic-rank-lower-bound}.
Exactness of the stable-part functor gives the same lower bound for
$\dim H^{n-1}(Y_0,\mathcal O_{Y_0})_{\mathrm{ss}}$.

For $d\geq\max\{1,r_\omega-1\}$, regularity and Serre duality give
$H^{n-1}(X,L^{-d})=H^1(X,\omega_X\otimes L^d)^\vee=0$.
The divisor sequence, including for the normal-crossings divisor,
then makes $h^{n-1}(Y,\mathcal O_Y)$ independent of $Y$ in the space of
properly intersecting degree-$d$ sections: the terminal map
$H^n(X,L^{-d})\to H^n(X,\mathcal O_X)$ is surjective because
$H^n(Y,\mathcal O_Y)=0$.  The universal family is
flat, and cohomology and base change therefore gives a vector bundle
with Frobenius.  If its rank is $h$, the stable rank at a perfect
geometric point is the rank of $F^h$; thus its maximal value is attained
on an open, and the generic value is at least the value at $Y_0$.
The smooth locus is dense and meets this open.

For a smooth $Y$, R1 identifies the stable ambient image with
$H^{n-1}(X,\mathcal O_X)_{\mathrm{ss}}$, of dimension $a$; in fact the
coherent restriction is injective in the present range of $d$.
Exactness of the stable-part functor subtracts $a$ and proves
\eqref{eq:generic-rank-lower-bound}.
For $d\geq2n$ one has
$\binom dn\geq(d/2)^n/n!$ and $\binom dm\leq d^{n-1}$ for $m<n$.
The right side of \eqref{eq:generic-rank-lower-bound} is therefore at
least $d^{n-1}(Dd/(2^nn!)-B)$, which is positive under
\eqref{eq:explicit-d0}.  Finally, Serre duality, Riemann--Roch, and
Serre vanishing give
$\dim E(Y,\mathcal O_Y)=Dd^n/n!+O_X(d^{n-1})$; together with the lower
bound this proves the asymptotic formula.
\end{proof}

\begin{corollary}[The eventual parenthesis of 8.5.8]
\label{cor:eventual-hodge-parenthesis}
For the explicit $d_0$ of \eqref{eq:explicit-d0}, every $d\geq d_0$
has the following properties.  Every smooth degree-$d$ section $Y$
satisfies
\[
 \dim E(Y,\mathcal O_Y)
 =h^0(X,\omega_X\otimes L^d)-h^0(X,\omega_X)>0.
\]
For every $Y$ in the maximal-stable-rank open $U_d$, the number of
$p$-adic unit roots in vanishing cohomology is positive.  Thus, on
this open and in this stated eventual range, unit-root nonvanishing
and the maximal Hodge-level condition $E(Y,\mathcal O_Y)\ne0$ are
equivalent.  No assumption on crystalline $p$-torsion of $X$ is needed.
\end{corollary}

\begin{proof}
The coherent divisor sequence and the vanishings used in the preceding
proof identify $E(Y,\mathcal O_Y)$ with
$\ker(H^n(X,L^{-d})\to H^n(X,\mathcal O_X))$, whose dimension is the
displayed, $Y$-independent difference by Serre duality.  It is positive
because it is at least the positive generic stable rank of
Theorem~\ref{thm:all-X-explicit-positive-rank}.
For $Y\in U_d$, Theorems~\ref{thm:A-exact85} and
\ref{thm:A-unit-count} identify the unit-root count with
\[
 f_{X,d}+\dim H^n_{\mathrm{cris}}(X/W)[p]_{\mathrm{ss}},
\]
which is positive.  Both properties therefore hold for every point
under these quantifiers.
\end{proof}

\par\smallskip\noindent\textbf{Scope of the manuscript's parenthesis.}
The supplied transcription of cote~4, scan~8, first specifies a
section for which the semisimple rank has its maximum value.  Scan~9,
continuing 8.5.8, then says ``pour tous ces $Y$'' and, equivalently,
``pour $Y$ g\'en\'erique''.  Corollary~\ref{cor:eventual-hodge-parenthesis}
answers this eventual generic assertion, including varieties with
$p$-torsion.  It does not assert that an arbitrary smooth section
outside $U_d$ has a unit root whenever its coherent vanishing space
is nonzero.  For example, a supersingular smooth plane cubic has
one-dimensional $H^1(\mathcal O)$ and no unit root.  That example
distinguishes the stronger pointwise assertion from the quantifiers
actually written in 8.5.8; it is not substituted for the requested
eventual assertion.

\par\smallskip\noindent\textbf{Relation to Koblitz.}
Koblitz (1975), I.\S5, Theorems~1--2, pp.~139--141, proves the
asymptotic invertibility estimate with an error of order $d^{n-1}$.
The argument above gives explicit constants in terms of the embedding
and the displayed cohomology groups.  Koblitz's I.\S6, Theorem~3,
p.~142, concerns generic complete intersections, not every hypersurface
section of an arbitrary fixed ambient variety.  His revised conjecture
in II.\S3, pp.~149--151, must not be cited as an established exact rank
formula.

\subsection{All residues for hypersurfaces in every projective space}

\begin{theorem}\label{thm:projective-all-residues}
Let $k$ be finite of characteristic $p$, let $n\geq2$, and let
$d\geq pn+1$.  Every nonempty open in the degree-$d$ hypersurface
parameter space of $\mathbf P^n_k$ contains, over a finite extension,
a smooth hypersurface with any prescribed point-count residue
modulo $p$.
\end{theorem}

\begin{proof}
Choose the integer representative $t\in\{0,\ldots,p-1\}$ of the
desired residue.  Over a sufficiently large finite extension
$\mathbf F_Q/k$, choose a general-position collection of $d$
geometric hyperplanes on which Frobenius has $t$ orbits of length $n$
and one orbit of length $d-nt>n$.  Existence follows by taking a
rational point in the open set of general-position configurations
in the product of Weil restrictions of the dual projective space
from the indicated extension degrees.  This parameter space is
geometrically integral, since over the algebraic closure it is a
product of projective spaces; the general-position open is nonempty.
Lang--Weil supplies rational points over arbitrarily large fields.
One may first form these Weil restrictions over a fixed finite base
field and then take extension degrees congruent to one modulo the
least common multiple of the prescribed orbit lengths; this preserves
their field factors while making the field arbitrarily large.
Distinctness of every geometric hyperplane forces the specified
orbit lengths to be exact.

The union $Y_0$ of these hyperplanes is defined over $\mathbf F_Q$.
A rational point lying on one hyperplane lies on its entire Frobenius
orbit.  An orbit of length $n$ has exactly one common intersection
point in $\mathbf P^n$, and that point is rational.  An orbit of
length greater than $n$ has empty intersection.  The $t$ points from
the shorter orbits are distinct, because any $n+1$ distinct
hyperplanes have empty intersection.  Therefore
\[
 \#Y_0(\mathbf F_Q)=t
\]
as an equality of integers.
The general-position hyperplane-union argument proves that
$H^{n-1}(Y_0,\mathcal O_{Y_0})$ has bijective Frobenius, of the full
generic rank $\binom{d-1}{n}$.  Thus $Y_0$ belongs to the same open
$W_d$ on which the entire Lang local system is defined.
The coherent trace congruence, together with the standard vanishing
of intermediate coherent cohomology for a projective hypersurface,
reads
\[
 \#Y(\mathbf F_{Q^e})\equiv
 1+(-1)^{n-1}\operatorname{Tr}(\operatorname{Frob}_{Q^e}^{\mathrm{geom}},V_Y)
 \pmod p.
\]
Apply finite Chebotarev on $W_d$ to the Frobenius class of $Y_0$ and
the given smooth open.  The resulting smooth hypersurface has the
same residue $t$.  No exact monodromy-group computation is required.
\end{proof}

\subsection{Current gates}
The odd-degree characteristic-two obstruction refutes the universal
full-$\mathrm{GL}$ expectation, even for geometrically connected
$X=\mathbf P^2$.  The all-degree characteristic-two theorem identifies
the exact geometric and arithmetic images: full $\mathrm{GL}_g(2)$
for even $d\ge4$, and the full stabilizer
$2^{g-1}\!:\mathrm{GL}_{g-1}(2)$ for odd $d\ge5$.
For plane curves, odd characteristic in degree five is the smallest
remaining monodromy case in this contribution.
For arbitrary $X$, the explicit asymptotic positivity bound above is
proved using a normal-crossings union of hyperplane sections.
An exact generic stable-rank formula for all $X$ and a determination
of all of their monodromy groups are not proved here.

% /root/monodromy, 19 November 2026.
% Complete mathematical construction; independently audited by /root.
\subsection{A smooth ambient surface whose relative Hasse--Witt map is always singular}
\label{subsec:relative-ordinarity-counterexample}
\status{Proved; independently audited by /root.  The finite Hopf-algebra checks
were executed; the final invariant-ring/Bertini search was not executed.}

The following construction refutes eventual bijectivity on
$E(Y,\mathcal O_Y)$ for arbitrary smooth ambient varieties.  The
ambient variety is smooth.  In particular, this does not reuse the
singular cones in Koblitz's counterexamples.

\begin{lemma}\label{lem:picard-finite-quotient}
Let $G$ be a finite commutative group scheme over a field $k$, and
let $Z\to X$ be a $G$-torsor, where $Z$ is a positive dimensional
projective complete intersection of dimension at least two and $X$
is smooth and geometrically connected.  If $G^D$ is connected, then
\[
 \operatorname{Pic}_{X/k}^0=G^D.
\]
\end{lemma}

\begin{proof}
The Koszul resolution gives $H^0(Z,\mathcal O_Z)=k$ and
$H^1(Z,\mathcal O_Z)=0$, compatibly with every extension of scalars;
smoothness or reducedness of $Z$ is unnecessary.  Therefore the
identity component of $\operatorname{Pic}_{Z/k}$ is trivial: its
Zariski tangent space is zero.  The kernel, as an fppf sheaf, of
\[
 \operatorname{Pic}_{X/k}\longrightarrow\operatorname{Pic}_{Z/k}
\]
is $G^D$.  Indeed, after trivializing the pullback of a line bundle,
its descent datum is a $G$-linearization of $\mathcal O_Z$.
All functions on $G\times Z$ come from $G$, because
$H^0(Z_R,\mathcal O)=R$ for every $k$-algebra $R$.  The descent
cocycle condition is precisely the condition that this unit be a
character $G_R\to\mathbf G_{m,R}$.  This proves the kernel assertion
without a surjectivity claim about the entire Picard scheme.
The connected group $\operatorname{Pic}_{X/k}^0$ maps to the trivial
identity component of $\operatorname{Pic}_{Z/k}$ and hence lies in
$G^D$.  Conversely the connected kernel $G^D$ lies in
$\operatorname{Pic}_{X/k}^0$.
\end{proof}

\par\smallskip\noindent\textbf{A specified finite group scheme and linear action.}
Work first over $\mathbf F_2$ and put
\[
 G=\operatorname{Spec}\mathbf F_2[a,b]/(a^2,b^2),\qquad
 \Delta(a)=a\otimes1+1\otimes a,\quad
 \Delta(b)=b\otimes1+1\otimes b+a\otimes a.
\]
Its group law is
$(a,b)(c,e)=(a+c,b+e+ac)$.  The basis $1,a,b,ab$ is the basis dual
to $1,t,t^2,t^3$ in the Hopf algebra of $\alpha_4$; the multiplication
rule is the divided-power rule
$f_if_j=\binom{i+j}{i}f_{i+j}$, truncated at index four.  Thus
$G=(\alpha_4)^D$.  Its Frobenius is zero, so
$V_{\alpha_4}=D(F_G)=0$, by Cartier duality for Verschiebung
\cite[Proposition~5.19]{EdixhovenVanDerGeerMoonenAV}.

Let $J$ be the direct sum of two nilpotent Jordan blocks of size
four on $\mathbf F_2^8$, with $J e_i=e_{i+1}$ in each block.
The formula
\begin{equation}\label{eq:alpha4-regular-action}
 (a,b)\longmapsto I+aJ+bJ^2+abJ^3
\end{equation}
is a faithful representation on $\mathbf A^8$ and acts on
$\mathbf P^7$.  Multiplication of these matrices gives exactly the
displayed group law.
The only subgroup of $G$ of order two is
$H=\{a=0\}\simeq\alpha_2$, and every nontrivial subgroup contains it.
One may see this directly from the restricted Lie algebra, with
basis $J,J^2$ and square operation $J^{[2]}=J^2$, $(J^2)^{[2]}=0$:
its only one-dimensional restricted subalgebra is $\langle J^2\rangle$.
Since $G$ has height one, its subgroup schemes correspond to these
restricted Lie subalgebras.  A projective line $\langle v\rangle$
is fixed by $H$ precisely when $J^2v=0$; a nilpotent operator has
no nonzero eigenvalue.  Consequently the nonfree locus is exactly
\[
 Z_{\rm bad}=\mathbf P(\ker J^2)=\mathbf P^3\subset\mathbf P^7.
\]
In particular, it has codimension four.

\par\smallskip\noindent\textbf{The surface and its fixed embedding.}
Let $Q=\mathbf P^7/G$ be the projective quotient and let
$Q^\circ=(\mathbf P^7\setminus Z_{\rm bad})/G$.  The quotient exists,
the map $\mathbf P^7\to Q$ is finite, and $Q^\circ$ is smooth;
these are Kothari's Theorem~3.3, p.~6
\cite{Kothari2026}.  The image of $Z_{\rm bad}$ has dimension three.
Fix a projective embedding $Q\hookrightarrow\mathbf P^M$ supplied
by a sufficiently divisible invariant Veronese ring.  Five general
hyperplanes of $\mathbf P^M$ cut out a smooth surface
\[
 X=Q\cap H_1\cap\cdots\cap H_5\subset Q^\circ.
\]
Their intersection avoids the three-dimensional bad locus and is
smooth by Bertini on $Q^\circ$.  The five pulled-back equations cut
out a surface $Z\subset\mathbf P^7$ which is a complete intersection;
the equality of the number of equations and the codimension makes
them a regular sequence in the polynomial ring.  The action on $Z$
is free and $Z/G=X$.  Geometric connectedness of $Z$, and then of
$X$, follows from the Koszul computation $H^0(Z,\mathcal O_Z)=k$.
This is also the particular dimension-two case of Kothari's
Proposition~3.4, pp.~6--7, and Corollary~3.7, p.~9; those statements
explicitly require smoothness of the quotient, not of $Z$.

This construction can be made unambiguous over a specified finite
field without leaving an unspecified ``general'' choice.  Use the
given ordered coordinates of $\mathbf P^7$ and their graded
lexicographic monomial order.  Compute the homogeneous invariant
ring by solving the coaction equations.  Every coordinate square is
invariant, because $a^2=b^2=0$, so this invariant ring is a finite
module over $\mathbf F_2[x_0^2,\ldots,x_7^2]$.  Its computation is
therefore a finite module-kernel calculation over an explicit
polynomial ring, rather than an unbounded search through degrees.
Take the least positive
Veronese degree which is generated in degree one.  Use its reduced
row-echelon monomial basis to embed $Q$.  Enumerate finite extensions
of $\mathbf F_2$ in increasing degree, each presented by the first
monic irreducible polynomial in lexicographic order, and then
enumerate ordered five-tuples of hyperplanes.  Retain the first
tuple for which the intersection is two dimensional, smooth, and
disjoint from the bad locus.  These conditions are decidable by
elimination and the Jacobian criterion.  The preceding nonempty
Bertini open proves termination.  Fix this output field $k$, this
surface $X$, and the resulting very ample bundle
$L=\mathcal O_{\mathbf P^M}(1)|_X$.
The description therefore defines one embedded surface by a finite
algorithm; the proof applies to every tuple passing these tests.
We do not claim to have numerically executed the invariant-ring and
Bertini search to identify the first tuple.

Lemma~\ref{lem:picard-finite-quotient} gives
\begin{equation}\label{eq:alpha4-picard-surface}
 \operatorname{Pic}_{X/k}^0=\alpha_4,\qquad
 h^1(X,\mathcal O_X)=1.
\end{equation}
Frobenius on this one-dimensional coherent cohomology is zero.
For a direct verification, the Cartan--Leray spectral sequence for
$Z\to X$ and $H^1(Z,\mathcal O_Z)=0$ identifies
$H^1(X,\mathcal O_X)$ with $H^1(G,k)$, the primitive elements of
$\mathcal O(G)$.  Those primitives are exactly $k a$; $b$ has the
additional term $a\otimes a$ in its coproduct, and $ab$ has the
terms $a\otimes b+b\otimes a$.  Absolute Frobenius sends $a$ to zero.

\begin{theorem}\label{thm:relative-ordinarity-fails-smooth}
For the fixed smooth projective geometrically connected surface
$(X,L)/k$ just constructed, let $r_\omega$ be a regularity bound for
its embedded dualizing sheaf, and put
$d_* =\max\{1,r_\omega-1\}$.
For every $d\geq d_*$ and every smooth section $Y\in|L^d|$, of
genus $g_Y$, one has
\[
 \dim E(Y,\mathcal O_Y)=g_Y-1,
 \qquad
 \dim E(Y,\mathcal O_Y)_{\mathrm{ss}}
       =f_Y\leq g_Y-2,
\]
where $f_Y$ is the $2$-rank of its Jacobian.  In particular the
Frobenius on $E(Y,\mathcal O_Y)$ is never bijective in this entire
range, including at the generic point.
\end{theorem}

\begin{proof}
Regularity and Serre duality give
$H^1(X,L^{-d})=H^1(X,\omega_X\otimes L^d)^\vee=0$.
Thus $H^1(X,\mathcal O_X)\to H^1(Y,\mathcal O_Y)$ is injective.
The induced homomorphism of group schemes
$\alpha_4=\operatorname{Pic}_{X/k}^0\to J_Y$ has zero Lie kernel.
Its finite kernel is therefore etale, but a subgroup of the local
group scheme $\alpha_4$ is local.  The kernel is trivial, and this
finite monomorphism is a closed immersion.

Naturality of Verschiebung and $V_{\alpha_4}=0$ put
$\alpha_4^{(2)}$ inside $\ker(V_{J_Y}:J_Y^{(2)}\to J_Y)$.
For an abelian variety of dimension $g_Y$ and $2$-rank $f_Y$, this
kernel has order $2^{g_Y}$ and its etale quotient has order $2^{f_Y}$;
its connected component therefore has order $2^{g_Y-f_Y}$.
Here the degree and point-count assertions are respectively
Propositions~5.20 and~5.22 of
\cite{EdixhovenVanDerGeerMoonenAV}: Frobenius is bijective on geometric
points, so $V\circ F=[2]$ identifies the geometric points of $\ker V$
with $J_Y[2](\bar k)$.
The connected subgroup $\alpha_4^{(2)}$ has order four, so
$g_Y-f_Y\geq2$.
Equation~\eqref{eq:alpha4-picard-surface} and coherent injectivity
give the dimension $g_Y-1$.  Since the ambient Frobenius is zero,
exactness of stable parts gives
$\dim E(Y,\mathcal O_Y)_{\mathrm{ss}}=f_Y$.
This proves the strict inequality and failure of bijectivity.
\end{proof}

\par\smallskip\noindent\textbf{What is and is not determined.}
This is a smooth counterexample to the universal eventual-relative-
ordinarity assertion.  It is compatible with the explicit positive
generic-rank theorem: for still larger degrees its generic rank is
positive and at least one less than its coherent dimension.  The
argument gives an exact structural obstruction and the displayed
upper bound; it does not assert an exact generic $2$-rank for this
surface.  The exact generic value for arbitrary $X$ remains a
separate determination.

\par\smallskip\noindent\textbf{Retained finite verification.}
The program \texttt{code/odd\_plane\_and\_alpha4\_checks.py} verifies
the Hopf group law, associativity, the matrix representation, its
divided-power multiplication table, the one-dimensional primitive
space and its zero Frobenius, and $\operatorname{rank}(J^2)=4$ on
the eight-dimensional representation.  The output is retained in
\texttt{outputs/odd\_plane\_and\_alpha4\_checks.json}.  These are
checks of the finite algebraic construction; existence and
smoothness of the surface are proved by the specified quotient and
Bertini construction, not inferred from a sample point calculation.

% Exact quintic case; independently audited by parent /root.
\subsection{The exact group for ordinary plane quintics in characteristic two}
\label{subsec:quintic-exact}

\begin{theorem}\label{thm:quintic-char2-monodromy}
For the universal ordinary smooth plane quintic in characteristic two,
the geometric and arithmetic monodromy groups on
$V=H^1_{\mathrm{et}}(C_{\bar\eta},\mathbf F_2)$ are both the full
stabilizer of the canonical nonzero covector:
\[
 G_{\mathrm{geom}}=G_{\mathrm{arith}}
 =\operatorname{Stab}_{\mathrm{GL}_6(\mathbf F_2)}(\ell)
 \simeq\mathbf F_2^5\rtimes\mathrm{GL}_5(\mathbf F_2).
\]
In particular their order is $319979520$, the representation is
reducible with invariant hyperplane of dimension five, and the
determinant character is trivial.
\end{theorem}
\begin{proof}

\par\smallskip\noindent\textbf{Theta-characteristic input, with a proof.}
For a smooth curve over an algebraically closed field of characteristic
two, let $B=F_*\mathcal O_C/\mathcal O_C$ for absolute Frobenius.
This is a line bundle with $B^{\otimes2}\simeq\omega_C$.
Locally a separating parameter $t$ gives the basis $1,t$ of
$F_*\mathcal O_C$ over $\mathcal O_C$; a change of separating parameter
$u=a^2+tb^2$ sends the generator of $B$ to $b$ times that generator,
whereas $du=b^2dt$.  This proves the claimed square relation.
The long exact sequence of
$0\to\mathcal O_C\to F_*\mathcal O_C\to B\to0$ shows
\[
 H^0(C,B)=\ker\bigl(F:H^1(C,\mathcal O_C)\to H^1(C,\mathcal O_C)\bigr).
\]
Thus $B$ is non-effective for an ordinary curve.
Every other theta characteristic is $B\otimes T$ for a nontrivial
two-torsion line bundle $T$.  Tensoring this sequence by $T$ and using
the projection formula and $F^*T=T^{\otimes2}\simeq\mathcal O_C$
gives an injection
\[
 H^0(C,\mathcal O_C)=H^0(C,F_*F^*T)
 \lhook\joinrel\longrightarrow H^0(C,B\otimes T),
\]
because $H^0(C,T)=0$.  Hence every theta characteristic other than $B$
is effective.  This is the statement recorded as St\"ohr--Voloch's
theorem in Ishitsuka--Ito, Theorem~3.1; Mumford (1971), p.~191,
describes the canonical characteristic-two theta characteristic.
For the geometric generic curve, the same reasoning may be performed
after perfecting its coefficient field.  Perfection does not change
the absolute Galois group, and $B$ is functorial under that group.

An ordinary curve of genus six has exactly $2^6=64$ geometric theta
characteristics.  On a plane quintic, $L=\mathcal O_C(1)$ is one of
them and $h^0(L)=3$, while $B$ is the unique non-effective one.
Under the identification of their set with the geometric points of
$J[2]$ by $T\mapsto L\otimes T$, the two distinguished points are
$0$ and $v=B\otimes L^{-1}\ne0$.  This is the dual representation
$V^\vee$, and $v$ is fixed by monodromy.

\par\smallskip\noindent\textbf{An irreducible incidence correspondence of contact conics.}
Let $S=\mathbf P(H^0(\mathbf P^2,\mathcal O(5)))$, of dimension twenty.
Let $I$ parametrize pairs $(C,Q)$ where $Q$ is a smooth plane conic
and the divisor cut out by $C$ on $Q$ is twice a divisor of degree five.
For a fixed conic $Q\simeq\mathbf P^1$, this means that the restriction
of its quintic equation, a binary form of degree ten, is a square.
The square forms constitute a vector subspace of dimension six inside
$H^0(Q,\mathcal O_Q(5))$, of dimension eleven.  The restriction map
from ternary quintics is surjective, with kernel
$Q\cdot H^0(\mathbf P^2,\mathcal O(3))$ of dimension ten.
Consequently the allowed quintic equations form a vector space of
dimension sixteen, or a projective space of dimension fifteen.
The smooth-conic parameter space has dimension five, so $\dim I=20$.
It is irreducible: fix $Q_0=(XZ+Y^2=0)$, take the irreducible product
$\mathrm{PGL}_3\times\mathbf P^{15}$ for its allowed quintics, and
map it surjectively to $I$ by changing projective coordinates.

The locus of pairs with a reduced reducible conic has dimension at
most nineteen.  Indeed, write $Q=L_1\cup L_2$ and $P=L_1\cap L_2$.
If $C$ is smooth and $C\cap Q$ is even, then $P\in C$: away from $P$
each of the two degree-five line sections must have even multiplicity,
so its multiplicity at $P$ is odd.  On each line the restriction is
therefore its linear equation for $P$ times the square of a binary
quadratic.  Each such restriction space has dimension three, and the
two restrictions automatically agree with value zero at $P$.
The combined restriction space has dimension six, so the quintic
equations again form a sixteen-dimensional vector space.
Pairs of distinct lines have dimension four, proving the bound
$4+15=19$.  Double-line conics cut exactly the divisors corresponding
to $L=\mathcal O_C(1)$.

Every effective theta characteristic $\theta$ and nonzero section $s$
gives a contact conic: $s^2\in H^0(C,\omega_C)=H^0(C,\mathcal O_C(2))$
lifts uniquely to a quadratic form on $\mathbf P^2$.
If that conic is a double line, then $\theta=L$.
Outside the dimension-at-most-nineteen image just described, each of
the other sixty-two theta characteristics therefore supplies a smooth
contact conic.  Thus $I\to S$ is dominant.  Since both dimensions are
twenty, its geometric generic fibre is zero dimensional.
It follows also that every such theta characteristic has $h^0=1$:
otherwise its projective space of sections would give a positive
dimensional family of conics in the same fibre.  Distinct theta
characteristics give distinct conics, because the half of an even
divisor on a smooth curve is unique.  The geometric generic fibre
therefore has exactly sixty-two points.

Irreducibility of $I$ implies transitivity of the geometric Galois
group on these sixty-two points.  This remains valid if the incidence
map has a purely inseparable part: the number of geometric generic
points is its separable degree, and a field's embeddings into a
separable closure form one orbit.  The Galois-equivariant
theta-characteristic interpretation proves that
$G_{\mathrm{geom}}$ on $V^\vee\setminus\{0,v\}$ is transitive.

\par\smallskip\noindent\textbf{The factor seven supplied by a quartic boundary.}
Take an ordinary smooth plane quartic $C_4$ and a line meeting it
transversely in four distinct points.  Their union $C_0$ is an ordinary
nodal quintic: the normalization sequence gives bijective Frobenius
on its $H^1(\mathcal O)$, of dimension six.  Over the parameter space
of such pairs, the etale normalization sequence gives a surjection
of local systems
\[
 H^1_{\mathrm{et}}(C_0,\mathbf F_2)
 \longrightarrow H^1_{\mathrm{et}}(C_4,\mathbf F_2).
\]
Forgetting the line is a smooth surjection with geometrically
connected fibres (nonempty opens of the dual plane), hence is
surjective on fundamental groups.  The quartic local system has
geometric monodromy $\mathrm{GL}_3(\mathbf F_2)$ by the already proved
quartic case.  Its order is divisible by seven, so the monodromy on
the nodal quintic family, and consequently the generic quintic group,
has order divisible by seven.  This specialization is legitimate
because the entire rank-six Lang system extends across $C_0$ on the
invertible-Hasse--Witt locus $W_5$.  The dense smooth open and $W_5$
have the same monodromy image for that extended finite etale system.

\par\smallskip\noindent\textbf{Finite group conclusion.}
Write $M=\operatorname{Stab}(v)\simeq2^5\!:\!\mathrm{GL}_5(2)$
in its action on $V^\vee$, and let $G\subset M$ be the geometric image.
Its transitive orbit of size sixty-two implies $31\mid|G|$; the
boundary argument gives $7\mid|G|$.
The complete maximal-subgroup list for $M$ is given in
Basheer--Moori (2021), Table~1, p.~3:
\[
\begin{array}{c|r|r}
\text{maximal subgroup}&\text{order}&\text{index in }M\\\hline
2_+^{1+8}:\mathrm{GL}_4(2)&10321920&31\\
2^{4+5}:\mathrm{GL}_4(2)&10321920&31\\
\mathrm{GL}_5(2)&9999360&32\\
2^{2+9}:(\mathrm{GL}_3(2)\times S_3)&2064384&155\\
2^9:(\mathrm{GL}_3(2)\times S_4)&2064384&155\\
(2^5:31):5&4960&64512.
\end{array}
\]
The first, second, fourth, and fifth orders are not divisible by
thirty-one.  The last is not divisible by seven.  The remaining
maximal subgroup is a complement $\mathrm{GL}_5(2)$; its action on
$V^\vee\setminus\{0,v\}$ has two orbits of length thirty-one, so it
cannot contain a group transitive on all sixty-two points.
No proper maximal subgroup can contain $G$, and therefore $G=M$.
The arithmetic image is contained in $M$ and contains $G$, hence
equals it.
\end{proof}

\par\smallskip\noindent\textbf{Independent elementary check of the complement issue.}
The same conclusion can be obtained from the maximal-subgroup table
for $\mathrm{GL}_5(2)$ in the ATLAS and the elementary fact
$H^1(\mathrm{GL}_m(2),\mathbf F_2^m)=0$ for $m\geq4$.
For completeness, let $c$ be a cocycle and put
$t_{ij}=1+E_{ij}$.  The relation $t_{ij}^2=1$ says that the $j$th
coordinate of $c(t_{ij})$ is zero.  Commuting $t_{ij}$ with
$t_{k\ell}$ for an auxiliary index $k$ distinct from $i,j,\ell$
shows that every coordinate except the $i$th is zero.
Thus $c(t_{ij})=a_{ij}e_i$.  The relation
$t_{ij}t_{jk}t_{ij}=t_{jk}t_{ik}$, for distinct indices,
gives $a_{ik}=a_{jk}$, so $a_{ij}=b_j$ depends only on $j$.
The coboundary associated with $b=(b_j)$ agrees with $c$ on all
elementary transvections, which generate the group.  Hence $c$ is
a coboundary.  The same result holds for the dual module by the
inverse-transpose automorphism.  Complements in $2^5:\mathrm{GL}_5(2)$
are therefore conjugate to the standard one, justifying the asserted
two-orbit description independently of the affine maximal-subgroup
classification.

% Exact characteristic-two plane sextic case; proof for independent audit.
\subsection{The exact group for ordinary plane sextics in characteristic two}
\label{subsec:sextic-exact}

\begin{theorem}\label{thm:sextic-full-monodromy}
For the universal ordinary smooth plane sextic in characteristic two,
the geometric and arithmetic monodromy images on
$V=H^1_{\mathrm{et}}(C_{\bar\eta},\mathbf F_2)$ are both
$\mathrm{GL}_{10}(\mathbf F_2)$.  Thus $V$ is irreducible and its
determinant character is trivial.
\end{theorem}

\begin{proof}
The generic genus is ten, and ordinarity gives $2^{10}=1024$
geometric theta characteristics.  The canonical characteristic $B_C$
is the unique non-effective one.  It is invariant under the
geometric Galois group, also when a purely inseparable extension is
needed to define it as a line bundle.  Identifying the set of theta
characteristics with $V^\vee$ using $B_C$ as origin, the remaining
1023 effective theta characteristics correspond to the nonzero
vectors.  We first prove that monodromy is transitive on them.

\par\smallskip\noindent\textbf{The contact reciprocity calculation.}
Let $C$ be a smooth plane curve of degree $d$, let $Q$ be a smooth
plane curve of degree $e=d-3$, and suppose $C\cap Q=2D$.
The half-intersection divisor defines the same zero-dimensional
subscheme $D$ on both curves: locally, in coordinates in which one
curve is $y=0$, the other has contact $y=u x^{2m}+\cdots$, and both
half divisors have ideal $(y,x^m)$.
Put
\[
 \theta_C=\mathcal O_C(D),\qquad
 \theta_Q=\mathcal O_Q(D-3H).
\]
These are theta characteristics.  The two ideal-sheaf sequences,
twisted by $e$, are
\[
 0\to\mathcal O_{\mathbf P^2}(-3)\to\mathcal I_D(e)
       \to\theta_C\to0,
 \qquad
 0\to\mathcal O_{\mathbf P^2}\to\mathcal I_D(e)
       \to\theta_Q\to0.
\]
Since $H^1(\mathbf P^2,\mathcal O(-3))=H^1(\mathbf P^2,\mathcal O)=0$,
they give the exact numerical relation
\begin{equation}\label{eq:contact-reciprocity}
 h^0(\theta_C)=1+h^0(\theta_Q).
\end{equation}

\par\smallskip\noindent\textbf{Smooth contact cubics.}
For $d=6$, every effective $\theta_C$ with section $s$ supplies a
contact cubic, because $s^2\in H^0(C,\omega_C)=H^0(C,\mathcal O_C(3))$
lifts uniquely to a ternary cubic.
If $Q$ is smooth, its theta characteristics are degree-zero line
bundles.  For an ordinary cubic, the only non-effective one is $B_Q$;
the other is $\mathcal O_Q$.  For a supersingular cubic, the sole
geometric theta characteristic is $\mathcal O_Q$.

The incidence corresponding to $\theta_Q=B_Q$ over the ordinary
smooth-cubic parameter space is irreducible of dimension 27.
To see this without a descent assumption on $B_Q$, pull the
coefficient space back along Frobenius (and work in an affine
coefficient chart).  The quotient of the relative Frobenius pushforward
by the structure sheaf supplies $B_Q$ on the Frobenius-twisted universal
curve.  This purely inseparable base change is irreducible and changes
no geometric Galois group.  The spaces
$H^0(Q,B_Q(3))$ form a vector bundle of rank nine, since its degree is
nine and the genus is one.  Squaring their sections gives the allowed
restrictions in $H^0(Q,\mathcal O_Q(6))$.
The restriction of ternary sextics is surjective, with kernel
$Q\cdot H^0(\mathbf P^2,\mathcal O(3))$ of dimension ten.
Thus the allowed sextics form, after this purely inseparable cover,
a nineteen-dimensional affine space, or a projective space of
dimension eighteen.  Smooth cubics have parameter dimension nine,
so the total is $9+18=27$, equal to the sextic parameter dimension.
The image of an irreducible space is irreducible.

The smooth-contact-cubic loci with effective $\theta_Q$ do not
dominate the sextic parameter space.  Their incidence dimensions are
at most 27 by the same degree-nine Riemann--Roch calculation; by
\eqref{eq:contact-reciprocity}, every such fibre over a sextic contains
the positive dimensional projective linear system of its
$\theta_C$, which has at least two sections.  Their images therefore
have dimension at most 26.  This remains true for the supersingular
cubic locus, whose base dimension is smaller.

\par\smallskip\noindent\textbf{Singular contact cubics.}
We show that these also occur only for sextics in a proper closed
subset of the 27-dimensional parameter space.
For every fixed cubic, the kernel of restriction of sextic forms has
dimension ten.  The following bounds are therefore bounds on the
dimension of the allowed restriction space, before adding ten and
projectivizing.

An irreducible nodal cubic has normalization $\mathbf P^1$ and two
points $P_+,P_-$ above the node.  Its stratum has dimension eight.
The pullback of $\mathcal O_Q(6)$ has degree eighteen.
Away from $P_+,P_-$, the orders of the restricted sextic are even.
At these two points the parities are equal.  If both are even, the
binary form is a square of a degree-nine form; the node gluing
condition gives one independent linear condition on that square root,
leaving dimension nine.  If both are odd, it is a fixed binary
quadratic with zeros $P_+,P_-$ times a square of degree eight, again
dimension nine; the gluing values are then both zero.  The incidence
dimension is consequently at most $8+9+10-1=26$.

An irreducible cuspidal cubic has normalization $\mathbf P^1$, one
point above its cusp, and parameter dimension at most seven.  All
orders of the pulled-back sextic are even, since its total degree is
eighteen and all orders away from that point are even.  It is therefore
a square of a degree-nine binary form, a space of dimension ten.
Any descent condition can only decrease the dimension.  The resulting
bound is $7+10+10-1=26$.  The dimensions eight and at most seven of
these strata follow directly by fixing a singular point: singularity
imposes three independent linear conditions on cubic forms, and a
repeated tangent cone imposes at least one further condition.

If the cubic is a transverse union of a line and a smooth conic, its
parameter dimension is seven.  On the two normalization components,
the sextic has degrees six and twelve.  Its odd orders, if any, must
occur at both of their two intersection points.  With no odd orders,
the square-root spaces have dimensions four and seven, subject to
two independent gluing conditions, giving nine.  With odd orders at
both points, extracting their fixed divisors leaves square-root
spaces of dimensions three and six, also giving nine.  The total
incidence dimension is at most 25.  If the line is tangent to the
conic, the cubic stratum has dimension at most six; each component's
restriction must have even order at the unique intersection point,
and the two square-root spaces satisfy at least their common-value
condition.  Their combined dimension is at most ten, again giving
at most 25.

For three distinct lines in a triangle, the parameter dimension is
six.  The possible odd-node patterns are either no nodes or all three
nodes.  In the first case the three four-dimensional square-root
spaces have three independent common-value conditions; in the second,
the three spaces have dimension three each.  Both bounds are nine,
giving total dimension at most 24.  For three concurrent lines the
parameter dimension is five.  Each line restriction has even order
at their sole common point, and the three square-root spaces satisfy
two common-value conditions.  Their combined dimension is at most
ten, again giving at most 24.

Finally a nonreduced cubic is $L^2M$, allowing $L=M$.
Its restriction to a smooth sextic can be even only if the line
section $M|C$ is even.  Sextics with such a line lie in dimension at
most 26: for a fixed line, binary sextics that are squares form a
four-dimensional subspace of a seven-dimensional restriction space;
the ternary restriction kernel has dimension 21.  Projectivizing and
allowing the line to vary gives $(4+21-1)+2=26$.
All possible singular cubics have now been covered.

Outside the proper loci just bounded, every one of the 1023 effective
theta characteristics of an ordinary sextic has a smooth contact
cubic, and this cubic corresponds to $B_Q$ by
\eqref{eq:contact-reciprocity}.  The irreducible incidence therefore
dominates the sextic space.  Its dimension is 27, so its generic
fibre is finite; every effective theta has $h^0=1$, and the fibre has
exactly 1023 geometric points.  As in the quintic case, irreducibility
implies geometric Galois transitivity on these points even if the
incidence has a purely inseparable part.  Thus $G_{\mathrm{geom}}$
is transitive on $V^\vee\setminus\{0\}$.

\par\smallskip\noindent\textbf{A boundary element that excludes the smaller transitive groups.}
Degenerate a sextic to an ordinary smooth quintic and a transverse
line, meeting in five distinct points.  This is an ordinary nodal
sextic.  Its rank-ten local system has a four-dimensional graph
submodule and the rank-six quintic local system as quotient.
The node permutation group is a subgroup of $S_5$; the graph
representation is its augmentation representation.
The exact quintic monodromy theorem supplies a quotient with order
divisible by seven and an element of order thirty-one acting on the
six-dimensional quotient as a one-dimensional identity plus an
irreducible five-dimensional summand.  Hence $7\mid|G_{\mathrm{geom}}|$.
Take the inverse image of the cyclic order-thirty-one subgroup in
the boundary monodromy group.  A Sylow thirty-one subgroup of this
inverse image maps onto the cyclic quotient.  There are no elements
of order $31^2$ in $\mathrm{GL}_{10}(2)$, since the multiplicative
order of $2$ modulo $31^2$ is 155.  Every nonidentity element of this
Sylow subgroup therefore has order thirty-one; choose one with
nonidentity quotient image and denote it by $g$.
It acts trivially on the graph submodule, since $31\nmid|S_5|$.
Its characteristic polynomial consequently has the form
\begin{equation}\label{eq:sextic-order31}
 \chi_g(T)=(T-1)^5 f(T),
\end{equation}
where $f$ is an irreducible polynomial of degree five whose roots
are primitive thirty-first roots of unity.
This polynomial is not reciprocal: inversion does not belong to
the orbit of squaring on primitive thirty-first roots, since
$\langle2\rangle\subset(\mathbf Z/31)^\times$ has odd order five.
As before, the specialization is inside the full-rank Lang locus,
so these are elements of the same generic monodromy image.

Apply Hering's classification of transitive linear groups in the
form of Churikov--Vasil'ev--Zvezdina (2022), Lemma~3.3, arXiv p.~6,
which was opened and checked.  In dimension ten over $\mathbf F_2$,
its alternatives, unless the group contains $\mathrm{SL}_{10}(2)$,
place the group in one of
\[
 N(\mathrm{Sp}_{10}(2)),\qquad
 \Gamma\mathrm L_5(4),\qquad
 \Gamma\mathrm L_2(32),\qquad
 \Gamma\mathrm L_1(1024).
\]
The exceptional alternatives have dimensions four or six or odd
characteristic and do not occur here; the $G_2$ alternative has
dimension divisible by six.
The symplectic normalizer preserves its unique invariant alternating
form up to a scalar, and over $\mathbf F_2$ the scalar is one.
Every one of its characteristic polynomials is reciprocal, contrary
to \eqref{eq:sextic-order31}.
An element of order thirty-one in $\Gamma\mathrm L_5(4)$ lies in
$\mathrm{GL}_5(4)$, so its fixed space has even dimension over
$\mathbf F_2$, contrary to the fixed dimension five in
\eqref{eq:sextic-order31}.
Finally, neither $\Gamma\mathrm L_2(32)$ nor
$\Gamma\mathrm L_1(1024)$ has order divisible by seven, which has
already been forced by the quintic quotient.
All proper alternatives are excluded.  Therefore
$G_{\mathrm{geom}}=\mathrm{SL}_{10}(2)=\mathrm{GL}_{10}(2)$.
The arithmetic image contains this group and is contained in it,
so it is the same group.
\end{proof}

% Characteristic-two plane septic case: separate pending independent audit.
\subsection{The exact group for ordinary plane septics in characteristic two}
\label{subsec:septic-exact}

\begin{theorem}\label{thm:septic-stabilizer-monodromy}
For the universal ordinary smooth plane septic in characteristic two,
the geometric and arithmetic monodromy images on
$V=H^1_{\mathrm{et}}(C_{\bar\eta},\mathbf F_2)$ are the full stabilizer
of the distinguished nonzero covector supplied by
$\mathcal O_C(2)$ and the canonical theta characteristic:
\[
 G_{\mathrm{geom}}=G_{\mathrm{arith}}
 =\operatorname{Stab}(\ell)
 \simeq 2^{14}\!:\mathrm{GL}_{14}(\mathbf F_2).
\]
In particular, the rank-fifteen representation is reducible, with its
distinguished invariant hyperplane, and its determinant is trivial.
\end{theorem}

\begin{proof}
We work first over an algebraic closure.  An ordinary genus-fifteen
curve has $2^{15}$ geometric theta characteristics.  Its canonical
theta characteristic $B_C$ is the unique non-effective one.  The
plane theta characteristic $L_C=\mathcal O_C(2)$ has six sections.
The class $v=L_C\otimes B_C^{-1}$ is the distinguished nonzero element
of $V^\vee$, fixed by arithmetic as well as geometric monodromy.
We prove transitivity on the other $2^{15}-2=32766$ elements of
$V^\vee$, then identify the whole stabilizer using a boundary subgroup.

\par\smallskip\noindent\textbf{Smooth contact quartics.}
Every nonzero section $s$ of an effective theta characteristic
$\theta_C$ supplies a unique quartic $Q$ with
$Q\cap C=2\operatorname{div}(s)$: the section $s^2$ belongs to
$H^0(C,\omega_C)=H^0(C,\mathcal O_C(4))$, and restriction from plane
quartics is an isomorphism.  If $Q$ is smooth, set
$D=\operatorname{div}(s)$ and
$\theta_Q=\mathcal O_Q(D-3H)$.  The contact reciprocity calculation
\eqref{eq:contact-reciprocity} gives
\[
 \theta_Q^2\simeq\omega_Q,\qquad
 h^0(\theta_C)=1+h^0(\theta_Q).
\]
For every smooth quartic and every theta characteristic on it,
$\theta_Q(3)$ has degree fourteen and exactly twelve sections, by
Riemann--Roch and Serre duality.  Restriction of septic forms to a
quartic is surjective and has kernel
$Q\cdot H^0(\mathbf P^2,\mathcal O(3))$, of dimension ten.

Consider first ordinary smooth quartics with $\theta_Q=B_Q$.
After a purely inseparable Frobenius base change, $B_Q$ is supplied
by the relative Frobenius sequence, so its twelve-dimensional spaces
of sections form a vector bundle.  Squaring their sections and lifting
through restriction gives an irreducible incidence of dimension
\[
 14+(12+10-1)=35,
\]
the dimension of the projective space of septics.  This incidence is
geometrically irreducible: the base is a nonempty open in the
geometrically irreducible quartic parameter space, the Frobenius
base change is a universal homeomorphism, and the section and lifting
spaces have geometrically irreducible affine fibres before
projectivizing.

For all smooth quartics and all theta characteristics, the incidence
dimension is at most thirty-five, since the scheme of roots is
finite over the quartic parameter space and the same twelve-section
calculation applies.  Those with effective $\theta_Q$ have images
of dimension at most thirty-four.  Indeed, reciprocity gives
$h^0(\theta_C)\ge2$, and the open part of
$\mathbf P H^0(C,\theta_C)$ near a smooth contact quartic is a
positive-dimensional fibre of this incidence over $C$.  This covers
the nonordinary smooth quartics as well, since all their geometric
theta characteristics are effective.  Thus a generic septic can
have a smooth contact quartic only through the canonical
non-effective theta characteristic of an ordinary quartic.

We next exclude singular contact quartics for the generic septic,
apart from the squares that give $L_C$.

\par\smallskip\noindent\textbf{Nonreduced quartics.}
A nonreduced quartic is either a square of a conic, or has the form
$Q=L^2R$ with $R$ a reduced conic.  A square gives exactly
$\theta_C=\mathcal O_C(2)$.  In the other case, evenness of $Q|C$
implies evenness of $R|C$.
For a fixed smooth conic, the restrictions of septics have binary
degree fourteen.  Those with even zero divisor are squares and form
an eight-dimensional space; the restriction kernel has dimension
$h^0(\mathbf P^2,\mathcal O(5))=21$.
Allowing the conic to vary gives dimension
$5+8+21-1=33$ for the image in the septic parameter space.
For a reduced pair of lines, each restricted binary form has odd
degree seven.  Its sole possible odd-order zero must be the node,
so each restriction is the fixed linear node factor times the
square of a cubic.  This leaves four parameters on each line;
both node values vanish, so their combined dimension is eight.
The conic parameter dimension is four, giving the smaller bound
$4+8+21-1=32$.  Consequently all non-square nonreduced contact
quartics occur over a proper subset of the septic parameter space.

\par\smallskip\noindent\textbf{Reduced reducible quartics.}
We give upper bounds that allow all coincidences and higher contacts
among components.  For a fixed reduced quartic, normalize each
component.  Away from the preimages of singular points, the divisor
of the restricted septic is even.  For each choice of the finite
odd-parity divisor $\epsilon$ on these preimages, its restriction
is a fixed section with divisor $\epsilon$ times a square.
The square-root line bundles form a finite set.  On a normalization
component of degree $a$ and genus $h$, the root has degree
$(7a-|\epsilon_i|)/2$ and, in the cases below, vanishing $H^1$;
hence its section dimension is
$(7a-|\epsilon_i|)/2+1-h$.
These degrees exceed $2h-2$: the only positive-genus component is
a smooth cubic, where the root degree is at least nine.
On a line the parity divisor has at most three points, and on a
conic at most four, so the genus-zero assertion is immediate.

For a smooth cubic and a line, the base dimension is eleven.
Both components have odd total intersection degree, so
$|\epsilon|\ge2$.  The sum of root dimensions is at most
$14+2-1-1=14$.  Adding the ten-dimensional restriction kernel and
projectivizing gives at most $11+14+10-1=34$.
If the cubic is irreducibly singular, its base dimension decreases
by at least one while the normalization genus decreases by one;
the same upper bound thirty-four follows.  Cases where the cubic
itself is reducible are included below.

For two smooth conics the base dimension is ten and the unrestricted
sum of root dimensions is sixteen.  If $\epsilon\ne0$, its total
degree is at least two, giving at most fifteen.  If
$\epsilon=0$, choose any intersection point of the two conics.
Descent of the septic restriction imposes equality of its two values,
and, after taking square roots over the algebraically closed field,
this is one independent linear condition on the root sections.
Independence follows from global generation on both normalization
components.  Thus the root dimension is again at most fifteen,
and the total bound is $10+15+10-1=34$.

For a smooth conic and two distinct lines the base dimension is
nine.  The two odd-degree line restrictions force
$|\epsilon|\ge2$, so the sum of root dimensions is at most
$14+3-1=16$.  The total is at most $9+16+10-1=34$.
For four distinct lines the base dimension is eight.  Each line
forces an odd parity divisor, so $|\epsilon|\ge4$ and the sum of
root dimensions is at most $14+4-2=16$; the total is at most
$8+16+10-1=33$.  No transversality assumption was used in these
upper bounds.

\par\smallskip\noindent\textbf{Integral singular quartics: normalization parameter spaces.}
An integral singular quartic has normalization of genus
$h\in\{0,1,2\}$.  We need two parameter bounds, and establish them
directly in characteristic two instead of using a Severi-variety
dimension assertion:
\begin{align*}
 \dim\{Q:\text{normalization genus }h\}&\le11+h,\\
 \dim\{Q:\text{normalization genus }h,
          \text{all singularities unibranch}\}&\le10+h.
\end{align*}
For $h=0$, a normalization map is given by a base-point-free
three-dimensional subspace of $H^0(\mathbf P^1,\mathcal O(4))$,
together with a projective frame.  The parameter dimension is
$\dim\operatorname{Gr}(3,5)+\dim\operatorname{PGL}_3
 -\dim\operatorname{PGL}_2=6+8-3=11$.
For $h=1$, temporarily mark an origin on the normalization.  Its
moduli have dimension one, its degree-four Picard variety dimension
one, the three-dimensional subspace of its four-dimensional section
space varies in dimension three, and a projective frame contributes
eight.  Forgetting the arbitrary origin subtracts one, giving
$1+1+3+8-1=12$.
For $h=2$, smooth-curve moduli have dimension three, the degree-four
Picard variety dimension two, and every degree-four line bundle has
three sections, giving $3+2+8=13$.
The dimension of genus-two moduli follows directly from curve
deformations: $H^0(T)=0$, $h^1(T)=3$ by Riemann--Roch, and $H^2(T)=0$.
The images under consideration are birational normalization maps,
so quotienting by reparametrization introduces no positive-dimensional
stabilizer.  These constructions cover all integral quartics of the
indicated normalization genus.

If all singularities are unibranch, the normalization map is
nonimmersive somewhere.  Otherwise at the preimage of each singular
point, one plane coordinate would be a uniformizer, making that
branch smooth.  For $h=0$ or $1$, and a degree-four line bundle $A$,
the map
$H^0(A)\to H^0(A|_{2P})$ is onto for every $P$, by Serre duality.
Among three-dimensional subspaces of $H^0(A)$, failure to have rank
two under this map has codimension two: for $h=1$ the rank-one locus
has dimension one in $\operatorname{Gr}(3,4)$ of dimension three;
for $h=0$ it has dimension four in
$\operatorname{Gr}(3,5)$ of dimension six.
Allowing $P$ to vary restores at most one dimension, so the
nonimmersive locus has codimension at least one in the parameter
space above.  This argument is invariant under reparametrization,
and the arbitrary elliptic origin remains a one-dimensional fibre.
For $h=2$, the system is the complete system $|A|$.
Failure to separate $2P$ is possible exactly when
$H^1(A(-2P))\ne0$, or $A\simeq\omega(2P)$.
As $P$ varies this is a subset of dimension at most one in the
two-dimensional degree-four Picard variety.  This remains a
dimension bound in characteristic two even if the map
$P\mapsto\omega(2P)$ is inseparable.  Thus the second bound follows
for all three genera.

Return to the restricted septic on such a quartic.  Its pullback
to the normalization has degree twenty-eight.  The odd-parity
divisor $\epsilon$ is supported above the finitely many singular
points.  There are at most six such branches, since the arithmetic
genus defect is at most three and the number of branches is at
most twice that defect.  A root after removing $\epsilon$ therefore
has degree at least eleven and vanishing $H^1$.  Its section
dimension is exactly
\[
 15-h-\frac{|\epsilon|}{2}.
\]
If $\epsilon\ne0$, then $|\epsilon|\ge2$; together with the
parameter bound $11+h$ this gives incidence dimension at most
$(11+h)+(14-h)+10-1=34$.
If $\epsilon=0$ and a singular point has two different branches,
descent imposes equality of the values at two distinct preimages.
The root line bundle separates these points, again by its degree,
so this is one independent linear condition after taking square
roots; the same bound thirty-four follows.
If all singularities are unibranch, the improved parameter bound
$10+h$ instead gives
$(10+h)+(15-h)+10-1=34$ without using any descent condition.
This excludes every integral singular contact quartic over a
generic septic.

\par\smallskip\noindent\textbf{Transitivity.}
Outside the proper closed subsets just bounded, every effective
theta characteristic other than $L_C$ has a smooth contact quartic,
and that contact quartic arises from its canonical non-effective
theta characteristic.  All $32766$ such theta characteristics
therefore occur in the geometrically irreducible canonical incidence
of dimension thirty-five.  This proves dominance.  Its generic
fibre is finite; in particular each of these theta characteristics
has exactly one section up to scalar, since otherwise its projective
section space would give a positive-dimensional fibre.  Consequently
the generic geometric fibre has precisely $32766$ points.
Irreducibility gives geometric Galois transitivity on these points;
a purely inseparable part of the incidence does not affect this
conclusion.  Hence $G_{\mathrm{geom}}$ is transitive on
$V^\vee\setminus\{0,v\}$.

\par\smallskip\noindent\textbf{The full quotient on the fourteen-dimensional space.}
Let $W=V^\vee/\langle v\rangle$, and let $H$ be the image on $W$.
Transitivity just proved implies transitivity of $H$ on
$W\setminus\{0\}$, of size $2^{14}-1$.
Degenerate to an ordinary sextic and a transverse line meeting it
in six points.  The rank-fifteen local system is still defined on
this nodal septic because Frobenius is bijective; its graph part
has rank five and its curve quotient has rank ten.  The projection
from this boundary family to the ordinary sextic parameter space
has geometrically connected fibres, so the quotient monodromy is
the full $\mathrm{GL}_{10}(2)$ established above.
In particular $11\mid|G_{\mathrm{geom}}|$.
The kernel of $\operatorname{Stab}(v)\to\mathrm{GL}(W)$ is a
two-group; therefore $11\mid|H|$.

Choose in the sextic quotient an order-thirty-one element acting
as a primitive five-dimensional block plus an identity of rank five.
An order-thirty-one lift exists by the Sylow argument used for
sextics: $\mathrm{GL}_{15}(2)$ has no element of order $31^2$, since
$\operatorname{ord}_{31^2}(2)=155$.
Its action on the graph part is trivial, since the node action
factors through $S_6$.  On $V^\vee$ it therefore has characteristic
polynomial $(T-1)^{10}f(T)$ for a primitive irreducible
degree-five factor $f$.  Quotienting by the fixed $v$ gives on $W$
\[
 (T-1)^9 f(T).
\]
This polynomial is not reciprocal, and the fixed-space dimension
is nine.

Hering's classification, in the opened form of
Churikov--Vasil'ev--Zvezdina (2022), Lemma~3.3, gives, in dimension
fourteen over $\mathbf F_2$, the alternatives
$\mathrm{GL}_{14}(2)$ or containment in
\[
 N(\mathrm{Sp}_{14}(2)),\quad
 \Gamma\mathrm L_7(4),\quad
 \Gamma\mathrm L_2(128),\quad
 \Gamma\mathrm L_1(16384).
\]
The symplectic alternative forces a reciprocal polynomial and is
excluded.  In $\Gamma\mathrm L_7(4)$ an order-thirty-one element
is $\mathbf F_4$-linear, so its fixed dimension over $\mathbf F_2$
is even, excluding nine.  The orders of the last two groups are
$1864254336$ and $229362$, respectively, neither divisible by eleven.
Thus $H=\mathrm{GL}_{14}(2)$.

\par\smallskip\noindent\textbf{The unipotent kernel and arithmetic equality.}
The intersection of $G_{\mathrm{geom}}$ with the kernel
$2^{14}$ of $\operatorname{Stab}(v)\to\mathrm{GL}(W)$ is an invariant
subspace for the full natural dual $\mathrm{GL}_{14}(2)$ module.
It is therefore zero or the whole kernel.  In the former case,
$G_{\mathrm{geom}}$ is a complement.  The elementary transvection
calculation $H^1(\mathrm{GL}_m(2),\mathbf F_2^m)=0$ for $m\ge4$,
proved in the quintic section (and equally valid for the dual),
shows that every such complement is conjugate to the standard one.
It has two orbits of size $2^{14}-1$ on
$V^\vee\setminus\{0,v\}$, contradicting transitivity.
Hence the whole kernel occurs, and the geometric group is the
full stabilizer.  Arithmetic monodromy also fixes $v$ and contains
the geometric group, so the two groups are equal.
\end{proof}

\par\smallskip\noindent\textbf{Finite certificate and scope.}
The retained program \texttt{sextic\_septic\_group\_checks.py}
checks the order of $2$ modulo $31^2$, the two distinct inverse
Frobenius orbits of primitive thirty-first roots, the semilinear
group orders and their missing prime divisors, and the two orbit
sizes used above.  It does not recompute Hering's classification
or replace the incidence argument by experimental Frobenius data.
This theorem settles degree seven in characteristic two only;
it does not assert the same stabilizer in every odd degree.

% All-degree characteristic-two contact argument; independent audit requested.
\subsection{Contact curves in every degree in characteristic two}
\label{subsec:all-degree-char2-contact}

The following argument removes the degree restriction from the
contact-incidence method.  Its global dimension input is
Christ--He--Tyomkin (2023), Proposition~2.7 (p.~8) and
Corollary~5.3 (p.~27), in the corrected publisher version of
\emph{On the Severi problem in arbitrary characteristic}.
Both statements apply over every algebraically closed field,
including characteristic two.  If a reduced plane curve of degree
$s$ has $c$ normalization components of genera $h_1,\ldots,h_c$,
their geometric genus convention is $1-c+\sum h_i$.  Consequently
the locus of such curves has dimension at most
\begin{equation}\label{eq:contact-severi-bound}
 3s+\sum_i h_i-c.
\end{equation}
If a family attains this dimension, its generic member is nodal.
The publisher's correction concerns the article history and missing
references; the opened corrected text contains these statements.

\begin{lemma}[The conductor inclusion for an even contact]
\label{lem:contact-conductor}
Let $C$ be a smooth plane curve of degree $d$ over an algebraically
closed field of characteristic two, and let $Q$ be a reduced plane
curve of degree $s<d$ with $Q\cap C=2D$ as divisors on $C$.
Write $\nu:\widetilde Q=\coprod Q_i\to Q$ for the normalization,
and write
\[
 \operatorname{div}(f|_{Q_i})=2A_i+\epsilon_i,
 \qquad M_i=\mathcal O_{Q_i}(A_i),
\]
where $f$ defines $C$ and $\epsilon_i$ is the reduced odd-order
divisor.  It is supported above the singular points of $Q$.
Set
\[
 J=\mathcal I_{D/Q}(s),\qquad
 L_C=\mathcal O_C(D),\qquad
 b=d-s.
\]
Then
\begin{equation}\label{eq:general-contact-reciprocity}
 h^0(C,L_C)=1+h^0(Q,J),
\end{equation}
and there is an inclusion
\begin{equation}\label{eq:conductor-section-inclusion}
 \nu_*\left(\bigoplus_i
       \omega_{Q_i}\otimes M_i^{-1}(3)\right)\hookrightarrow J.
\end{equation}
In particular, with $t=\sum_i h^1(Q_i,M_i)$,
\begin{equation}\label{eq:contact-h1-cancellation}
 h^0(Q,J)\ge t,
 \qquad t>0\Longrightarrow h^0(Q,J)\ge t+1.
\end{equation}
\end{lemma}

\begin{proof}
The same subscheme $D\subset\mathbf P^2$ is used on both curves.
Its ideal-sheaf sequences, twisted by $s$, are
\[
 0\to\mathcal O_{\mathbf P^2}(s-d)
  \to\mathcal I_D(s)\to L_C\to0,
 \qquad
 0\to\mathcal O_{\mathbf P^2}
  \to\mathcal I_D(s)\to J\to0.
\]
Here $\mathcal O_C(s)(-D)\simeq\mathcal O_C(D)$ because
$2D=Q|C$.  Since $s<d$ and the first cohomology of every line
bundle on $\mathbf P^2$ vanishes, these sequences prove
\eqref{eq:general-contact-reciprocity}.

We prove the inclusion locally and then express it globally.
At a point of $D$, complete the local ring of the plane.  Since
$C$ is smooth, choose its equation $f$ as one parameter and let
$b_0$ cut out the half divisor on $C$.  After multiplying the local
equation of $Q$ by a unit, that equation is
$b_0^2-fa_0$.  In the reduced one-dimensional ring
$A=\widehat{\mathcal O}_{Q,P}$, the ideal of $D$ is
$I=(f,b_0)$ and $f$ is a non-zero-divisor.
Reduction modulo $f$ gives
$A/(f)\simeq k[[x]]/(b_0^2)$, where the annihilator of $b_0$ is
$(b_0)$.  Hence
\[
 I^*:=\operatorname{Hom}_A(I,A)
 =\{z:zI\subset A\}=f^{-1}I,
 \qquad I^{**}=I.
\]
These equalities can be read directly in the total quotient ring:
$zf\in A$ and $zb_0\in A$ say that $zf$ belongs to
$(f):(b_0)=I$.

Let $v_i$ be the valuation on a normalization branch.  Put
\[
 r_i=v_i(f),\qquad j_i=\min\{v_i(f),v_i(b_0)\},
 \qquad\gamma_i=2j_i-r_i.
\]
The relation $b_0^2=fa_0$ gives $\gamma_i\ge0$; moreover
$\gamma_i\equiv r_i\pmod2$.  Thus
$\gamma_i=\epsilon_i(P_i)+2k_i$ for an integer $k_i\ge0$.
Write $c_i$ for the conductor exponent on this branch, so that
the conductor ideal on the normalization is generated there by
a parameter to the power $c_i$.
If a tuple $z$ of branch functions satisfies
\[
 v_i(z)\ge c_i+r_i-j_i\quad\hbox{for every branch},
\]
then its product with every member of $I^*=I/f$ lies in the
conductor, hence in $A$.  It follows that $z\in I^{**}=I$.
No separate assumption on a tame singularity, on its number of
branches, or on its analytic type is used here.

Let $\Delta$ be the conductor divisor on $\widetilde Q$, let
$\Gamma=\sum\gamma_iP_i$, and let $K_0=\sum k_iP_i$.
If $N=\nu^*J/\text{torsion}$, the valuation calculation gives
\[
 N|_{Q_i}\simeq M_i(-b)(-K_0|_{Q_i}),\qquad
 \nu_*N(\Gamma-\Delta)\subset J.
\]
Indeed the normalized ideal has vanishing order $j_i$, and
$\gamma_i=2j_i-r_i$ is precisely the adjustment from its square
to the divisor of $f$.

For a reduced plane curve, finite duality (Stacks Project,
Remark~48.20.12, Tag~0AX4, opened and checked) and adjunction give
\[
 \omega_{\widetilde Q}
 \simeq\nu^*\omega_Q(-\Delta)
 =\mathcal O_{\widetilde Q}(s-3)(-\Delta).
\]
For completeness, the conductor occurs because
$\operatorname{Hom}_{\mathcal O_Q}
(\nu_*\mathcal O_{\widetilde Q},\mathcal O_Q)$ is the conductor
ideal; finite duality identifies its tensor product with
$\omega_Q$ with $\nu_*\omega_{\widetilde Q}$.
Using $M_i^2\simeq\mathcal O_{Q_i}(d)(-\epsilon_i)$ now yields
\[
 \omega_{Q_i}\otimes M_i^{-1}(3)
 \simeq M_i(-b)(\epsilon_i-\Delta_i)
 =N|_{Q_i}(K_{0,i}+\epsilon_i-\Delta_i).
\]
Since $\Gamma=\epsilon+2K_0$, this line bundle is a subsheaf
of $N(\Gamma-\Delta)$, proving
\eqref{eq:conductor-section-inclusion}.

By Serre duality, $h^1(M_i)=h^0(\omega_{Q_i}M_i^{-1})$.
Twisting by the positive line bundle $\mathcal O_{Q_i}(3)$
does not decrease this dimension.  If the initial space is nonzero,
the increase is strict: choose a point where some section does not
vanish, and choose a plane cubic through its image whose pullback
is a nonzero divisor.  Multiplication by that cubic has image
consisting of sections vanishing there, whereas multiplying the
chosen initial section by a cubic not through the point gives a
section outside the image.  Summing over the components proves
\eqref{eq:contact-h1-cancellation}.
\end{proof}

\begin{lemma}[No lower-degree reduced contact for a generic curve]
\label{lem:no-low-degree-contact}
Fix $d\ge4$ and $0<s<d-3$ in characteristic two.
A generic smooth plane curve $C$ of degree $d$ has no reduced
plane curve $Q$ of degree $s$ with $Q|C$ even.
\end{lemma}

\begin{proof}
Let $N_d=\binom{d+2}{2}-1$.
For fixed normalization genera, component degrees, and odd-parity
divisors $\epsilon_i$, the line bundles
$M_i^2=\mathcal O_{Q_i}(d)(-\epsilon_i)$ have finite choices.
They can be parameterized after a finite base change of the
singularity stratum; inseparable base changes do not affect
dimension.  The normalized restrictions of $f$ are squares of
sections of these $M_i$, times the fixed odd-parity factors.
Ignoring descent conditions only increases the dimension.
Riemann--Roch gives
\[
 \sum_i h^0(M_i)
 =\frac{ds}{2}-\frac{|\epsilon|}{2}
   +c-\sum_i h_i+t,
 \qquad t=\sum_i h^1(M_i).
\]
The kernel of restricting degree-$d$ forms to $Q$ has dimension
$\binom{d-s+2}{2}$.  Combining this with
\eqref{eq:contact-severi-bound}, and projectivizing, gives the
incidence bound
\begin{equation}\label{eq:general-contact-dimension}
 \dim\mathcal I_{d,s}
 \le N_d-\frac{s(d-s-3)}2-\frac{|\epsilon|}2+t
\end{equation}
on each stratum with fixed $t$.

Here is the precise fibre argument used with this bound.
Suppose a component of the even-contact incidence dominates the
ordinary smooth degree-$d$ parameter space.  Over its generic
curve, roots $L_C^2=\mathcal O_C(s)$ form a finite scheme.
For each geometric root the contact forms are the projective
space of its nonzero sections, followed by squaring.
Restriction of degree-$s$ plane forms to $C$ is an isomorphism,
so no further choices of $Q$ occur.  The reduced-contact condition
is open in this projective section space.  The generic fibre of
our component therefore has dimension
$h^0(L_C)-1=h^0(J)$.
Its generic point belongs to a single stratum of the finite
stratification by the data above, and an open subset of the same
generic fibre remains in that stratum.  Thus the fibre dimension
may be subtracted in \eqref{eq:general-contact-dimension}; we are
not subtracting a family of contact curves from a special stratum
that it need not preserve.

The conductor lemma gives $h^0(J)\ge t$.  Dominance would imply
\[
 N_d\le N_d-\frac{s(d-s-3)}2-\frac{|\epsilon|}2<N_d,
\]
a contradiction.  There are finitely many relevant component
degrees, genera, parity choices in finite covers, and cohomology
dimensions.  Thus their images have proper closure, proving the
claim for a generic curve.
\end{proof}

\begin{theorem}[The theta orbits of generic plane curves in characteristic two]
\label{thm:all-char2-theta-orbits}
Let $d\ge4$, $g=(d-1)(d-2)/2$, and let $C$ be the geometric
generic ordinary smooth plane curve of degree $d$ in characteristic
two.
If $d$ is even, every theta characteristic other than the canonical
one $B_C$ has exactly one section, and geometric monodromy is
transitive on these $2^g-1$ characteristics.
If $d$ is odd, every theta characteristic other than $B_C$ and
$\mathcal O_C((d-3)/2)$ has exactly one section, and geometric
monodromy is transitive on these $2^g-2$ characteristics.
Equivalently, on $V^\vee$ geometric monodromy is transitive on
all nonzero vectors in even degree, and on
$V^\vee\setminus\{0,v\}$ in odd degree, where $v$ is the
distinguished nonzero invariant vector.
\end{theorem}

\begin{proof}
Put $e=d-3$.  A section of an effective theta characteristic
$\theta_C$ gives a contact form $Q$ of degree $e$, and conversely
such a form gives an effective theta characteristic and its
section up to scalar.  Work throughout over the ordinary smooth
degree-$d$ locus, so theta characteristics have finitely many
geometric choices, with constant count $2^g$.  Passing to a
purely inseparable cover if necessary changes neither dimensions
nor geometric Galois orbits.

\par\smallskip\noindent\textbf{The unique dominant smooth-contact incidence.}
For a smooth $Q$, contact reciprocity gives
$h^0(\theta_C)=1+h^0(\theta_Q)$, where
$\theta_Q=\mathcal O_Q(D-3)$.
Riemann--Roch gives
\[
 h^0(\theta_Q(3))=3e+h^0(\theta_Q(-3)).
\]
Thus the incidence dimension, for a stratum with these dimensions
constant, is at most $N_d+h^0(\theta_Q(-3))$.
If $\theta_Q$ is effective, then
$h^0(\theta_Q(-3))<h^0(\theta_Q)$: multiplication by a suitable
cubic gives a proper subspace, as in the final paragraph of the
conductor lemma.  The fibre dimension $h^0(\theta_C)-1$
is $h^0(\theta_Q)$, so such an incidence cannot dominate.
For an ordinary smooth $Q$, its canonical $B_Q$ is the unique
non-effective theta characteristic.  It has
$h^0(B_Q(3))=3e$.  For a nonordinary smooth $Q$, all its geometric
theta characteristics are effective.  Consequently the only
possible dominant smooth-contact incidence is the canonical one
over ordinary smooth $Q$.

That canonical incidence is geometrically irreducible of dimension
$N_d$: the ordinary degree-$e$ base is geometrically irreducible;
after Frobenius base change the relative canonical theta
characteristic exists; its $3e$ sections form a vector bundle;
and lifts of their squares through restriction form affine spaces
with the constant ten-dimensional kernel
$Q\cdot H^0(\mathbf P^2,\mathcal O(3))$.
Projectivization gives total dimension
$\binom{e+2}{2}-1+3e+10-1=N_d$.
For $e=1,2$ the genus-zero canonical theta characteristic is used,
and the same computation applies.

\par\smallskip\noindent\textbf{Reduced singular contact forms cannot dominate.}
Apply the notation and incidence bound of the preceding lemmas
with $s=e$.  A hypothetical dominating component has dimension
at most $N_d-|\epsilon|/2+t$, and generic fibre dimension
$h^0(J)$.  If $t>0$, the strict part of
\eqref{eq:contact-h1-cancellation} gives a contradiction.
If $|\epsilon|>0$, its total degree is even and at least two,
so the non-strict inequality also gives a contradiction.
If $h^0(J)>0$ while $t=|\epsilon|=0$, the fibre dimension again
gives a contradiction.  Only
\[
 t=0,\qquad \epsilon=0,\qquad h^0(J)=0
\]
could remain for a dominant component.

In this case its image in the degree-$e$ curve parameter space
must attain the maximal Severi dimension
$3e+\sum h_i-c$; otherwise the incidence bound loses at least
one further dimension.  By Christ--He--Tyomkin Corollary~5.3,
the generic $Q$ is therefore nodal.  Since it is singular, choose
one node.  The conductor inclusion and $h^0(J)=0$ imply
\[
 H^0(Q_i,\omega_{Q_i}M_i^{-1}(3))=0
 \quad\hbox{on every component}.
\]
If the two preimages $P_+,P_-$ of the node lie on one component,
there exists a plane cubic whose pullback contains both points.
It follows that
$H^0(\omega M_i^{-1}(P_++P_-))=0$, hence
$H^1(M_i(-P_+-P_-))=0$.  The evaluation map at the two points
is therefore onto.  If they lie on different components, the
same argument with one point proves surjectivity of each
one-point evaluation.  In either case, equality of the two values
needed for descent of the squared restriction imposes one
nonzero linear condition on the root sections after taking square
roots of its scalar coefficients.  This lowers the incidence
dimension by at least one, the final contradiction.

\par\smallskip\noindent\textbf{Nonreduced contact forms.}
Every nonreduced form of degree $e$ has a unique expression up to
scalars $Q=R^2S$, with $S$ reduced and $\deg S=s<e$.
If $s>0$, evenness of $Q|C$ implies evenness of $S|C$, excluded
for a generic $C$ by Lemma~\ref{lem:no-low-degree-contact}.
If $s=0$, then $e$ is even, and the half divisor of $Q|C$ is
the divisor of $R|C$.  It gives exactly the plane theta
characteristic $\mathcal O_C(e/2)$.  This case occurs precisely
when $d$ is odd and is the sole nonreduced contact case retained.

\par\smallskip\noindent\textbf{Counting the generic fibre and obtaining the orbit.}
Every noncanonical theta characteristic is effective.
All except the plane characteristic in odd degree must therefore
occur in the geometrically irreducible canonical smooth-contact
incidence.  This proves its dominance.  Since it has dimension
$N_d$, its generic fibre is finite, so each of these theta
characteristics has exactly one section up to scalar.
The generic geometric fibre has respectively $2^g-1$ or
$2^g-2$ points.  Irreducibility makes geometric Galois monodromy
transitive on them; purely inseparable factors of the incidence
do not change this orbit statement.  Translating theta
characteristics by the canonical $B_C$ identifies their set
with $V^\vee$, proving the asserted vector orbits.
\end{proof}


\subsection{Exact characteristic-two plane monodromy in every degree}
\label{subsec:char2-uniform}

\begin{lemma}[Two elementary group-theoretic inputs]
\label{lem:char2-group-inputs}
For every integer $a\ge8$, there is an odd prime $\ell_a$ with
$\operatorname{ord}_{\ell_a}(2)=a$.
For $N\ge4$, the first cohomology of $\mathrm{GL}_N(\mathbf F_2)$
with coefficients in its natural module or its dual is zero.
\end{lemma}
\begin{proof}
We include a proof of the needed special case of the primitive-divisor
theorem. If an odd prime $\ell$ divides $\Phi_a(2)$ without having
order $a$, the cyclotomic factorization and the lifting-the-exponent
identity show that
$a=\operatorname{ord}_{\ell}(2)\ell^j$ with $j\ge1$, and that
$v_\ell(\Phi_a(2))=1$. Indeed, substitute
$v_\ell(2^{fm}-1)=v_\ell(2^f-1)+v_\ell(m)$, where
$f=\operatorname{ord}_\ell(2)$, in the M\"obius formula for
$v_\ell(\Phi_a(2))$; the sum is zero unless $a/f$ is a power of
$\ell$, when it is one. Since $f\mid\ell-1$, such an $\ell$ is the
largest prime factor $P$ of $a$. Thus, if no primitive prime occurs,
$\Phi_a(2)$ is either $1$ or $P$.
On the other hand the cyclotomic product gives
\[
 \Phi_a(2)>2^{\varphi(a)}\prod_{j\ge1}(1-2^{-j})
 >2^{\varphi(a)-2}.
\]
For the last bound, the first three factors have product $21/64$,
and the remaining product is at least
$1-\sum_{j\ge4}2^{-j}=7/8$; their product exceeds $1/4$.
If $P\ge7$, then $\varphi(a)\ge P-1$ and $2^{P-3}>P$.
If $P\le3$ and $a\ge8$, then $\varphi(a)\ge4$ and the same
conclusion holds. If $P=5$, either $\varphi(a)\ge8$ or $a=10$;
the exceptional calculation is $\Phi_{10}(2)=11$.
This contradiction proves the first assertion. The cyclotomic
valuation calculation also agrees with
\cite[Lemma~3.1(b)--(c), pp.~125--126]{GlasbyEtAl2017}, whose
statement and proof were opened.

For the second assertion let $z$ be a $1$-cocycle with values in
$\mathbf F_2^N$ and put $t_{ij}=1+E_{ij}$.
The relation $t_{ij}^2=1$ forces the $j$th coordinate of
$z(t_{ij})$ to vanish. For $l\ne i,j$, choose
$k\notin\{i,j,l\}$. Comparing the cocycle on the commuting pair
$t_{ij},t_{kl}$ shows that its $l$th coordinate also vanishes.
Thus $z(t_{ij})=a_{ij}e_i$.
For distinct $i,j,k$, apply the cocycle to
$[t_{ij},t_{jk}]=t_{ik}$; it gives $a_{ik}=a_{jk}$.
Consequently $a_{ij}=b_j$ for a single vector $b$, and
$z(t_{ij})=(t_{ij}-1)b$. Since these transvections generate the
group, $z$ is a coboundary. Composing with inverse transpose
gives the dual-module assertion.
\end{proof}

\begin{theorem}[All plane degrees in characteristic two]
\label{thm:plane-char2-all-groups}
Let $k$ be any perfect field of characteristic two, let $U_d$ be
the ordinary smooth locus in the projective space of degree-$d$
plane forms, and choose a geometric generic base point $\bar\eta$.
Write
\[
 g=(d-1)(d-2)/2,\qquad
 V=H^1_{\mathrm{et}}(C_{\bar\eta},\mathbf F_2).
\]
For $d\ge4$, the images of $\pi_1(U_d\otimes_k\bar k,\bar\eta)$
and $\pi_1(U_d,\bar\eta)$ on $V$ are equal and are
\[
 \begin{cases}
 \mathrm{GL}_g(\mathbf F_2),&d\text{ even},\\
 \operatorname{Stab}(\ell)
   \simeq(\mathbf F_2)^{g-1}\rtimes
                  \mathrm{GL}_{g-1}(\mathbf F_2),&d\text{ odd},
 \end{cases}
\]
where $\ell\ne0$ is the canonical covector determined by
$\mathcal O_C((d-3)/2)\otimes B_C^{-1}$ in odd degree.
The nonzero representation is irreducible exactly in even degree;
its determinant is always trivial. In degrees one and two $V=0$;
in degree three it has rank one and trivial, hence irreducible,
monodromy. Thus there is no threshold after which all plane
degrees have full linear or irreducible monodromy.
\end{theorem}
\begin{proof}
The degrees at most four, the quintic stabilizer, and the sextic
full group have been proved above. We induct on $d\ge7$.
Work first over $\bar k$ and on the dual space $W_0=V^\vee$.
Theorem~\ref{thm:all-char2-theta-orbits} proves transitivity on
$W_0\setminus\{0\}$ when $d$ is even and on
$W_0\setminus\{0,\ell\}$ when $d$ is odd.
In the odd case put $W=W_0/\langle\ell\rangle$; in the even
case put $W=W_0$. The induced group on $W$ is therefore transitive
on its nonzero vectors. Its dimension is
\[
 N=\begin{cases}g,&d\text{ even},\\g-1,&d\text{ odd}.\end{cases}
\]

We specify the boundary subgroup used in the induction. The
universal plane-curve coherent $H^1$ is a vector bundle of rank
$g$, including singular fibres. The locus where its Frobenius is
bijective is open, and the Lang construction gives an etale
$\mathbf F_2$-local system on this whole locus. It contains the
nodal curves $C'\cup L$, where $C'$ is ordinary smooth of degree
$d-1$ and $L$ meets it transversely in $d-1$ distinct points.
The smooth ordinary locus is dense in this normal parameter
space. Restriction of connected finite etale covers to a dense
open remains connected, so its fundamental group maps onto the
fundamental group of the larger locus. Boundary monodromy thus
lies in the group being determined.

The normalization sequence for $C'\cup L$ gives
\[
 0\longrightarrow H^1(\Gamma,\mathbf F_2)
 \longrightarrow H^1_{\mathrm{et}}(C'\cup L,\mathbf F_2)
 \longrightarrow H^1_{\mathrm{et}}(C',\mathbf F_2)
 \longrightarrow0.
\]
Here the graph space has dimension $d-2$ and its monodromy factors
through the permutation group $S_{d-1}$ of the nodes. The
projection from the parameter space of transverse pairs $(C',L)$
to that of $C'$ is smooth, surjective, and has geometrically
connected fibres, namely nonempty opens of the dual plane.
It therefore surjects on fundamental groups and supplies the
entire already determined monodromy of $C'$.

Write $r=(d-2)(d-3)/2$. If $d$ is even, the preceding odd-degree
group has a quotient $\mathrm{GL}_{r-1}(2)$; the kernel of this
quotient in the boundary group has prime divisors only $2$ and
those dividing $(d-1)!$. If $d$ is odd, the dual boundary
sequence has an invariant subspace of dimension $r$ with full
$\mathrm{GL}_r(2)$ action. The nonzero invariant vector $\ell$
cannot lie in that subspace, since it has no nonzero invariant
vector. It remains nonzero at the boundary: an invariant section
of a finite etale sheaf on the dense open extends uniquely over
the normal larger locus, and the zero and nonzero sections are
disjoint. Thus the rank-$r$ subspace injects into $W$, with quotient
a subquotient of the graph permutation space. In this case also,
the kernel of the boundary quotient onto $\mathrm{GL}_r(2)$ has
prime divisors only $2$ and those dividing $(d-1)!$. These kernel
claims follow by writing the representation in block triangular
form: its off-diagonal part is a $2$-group, the preceding odd-degree
unipotent radical is a $2$-group, and its remaining diagonal
action factors through $S_{d-1}$.

Set $a=\lfloor N/2\rfloor+1$. The elementary inequalities
\[
 8\le a<a+1<N-1,\qquad a>d-1,\qquad
 a+1\le\begin{cases}r-1,&d\text{ even},\\r,&d\text{ odd}
 \end{cases}
\]
hold for even $d\ge8$ and odd $d\ge7$. For example the two initial
cases have $(N,r,a)=(21,15,11)$ and $(14,10,8)$.
For the upper bound, using $\lfloor N/2\rfloor\le N/2$
reduces it to $g\ge2d+2$ in even degree and $g\ge2d-1$ in odd
degree; both hold at the indicated initial degree and their
left-minus-right differences increase. The other inequalities
follow directly from the same genus formula.
For $j=a,a+1$, choose the prime $\ell_j$ in
Lemma~\ref{lem:char2-group-inputs}. We have $\ell_j>j>d-1$.
Also $\ell_j\nmid N$: otherwise $\ell_j>N/2$ forces $N=\ell_j$,
whereas $j\mid N-1$ and $N/2<j<N-1$, an impossibility.

In the boundary general linear quotient choose an element of
order $\ell_j$ with one irreducible $j$-dimensional block and
identity on its complement. The kernel has order prime to
$\ell_j$, so a Sylow subgroup of its inverse image supplies an
element of order $\ell_j$ with that image. It acts trivially on
the graph quotient, because $\ell_j\nmid(d-1)!$.
Semisimplicity of this odd-order action shows that its
characteristic polynomial on $W$ is
\[
                    (T-1)^{N-j}f_j(T),
 \qquad f_j\text{ irreducible of degree }j.
\]

Apply the opened classification in
\cite[Lemma~3.3, p.~6]{ChurikovVasilevZvezdina2022}.
Every proper extension-field alternative lies inside
$\Gamma\mathrm L_{N/b}(2^b)$ for some divisor $b>1$ of $N$.
Our order-$\ell_j$ elements belong to its linear subgroup,
since $\ell_j\nmid b$. Their fixed-space dimensions are therefore
divisible by $b$. The consecutive dimensions $N-a$ and $N-a-1$
force $b=1$, excluding all these alternatives, including the
one-dimensional semilinear case. Over the original field the
remaining symplectic alternative lies in $\mathrm{Sp}_N(2)$:
its normalizer preserves the unique invariant alternating form,
and the only nonzero scalar is one. One of $a,a+1$ is odd.
For this $j$, the irreducible factor $f_j$ is not reciprocal,
since inversion is not in the odd-order subgroup generated by
$2$ modulo $\ell_j$. Its characteristic polynomial contradicts
preservation of a nondegenerate alternating form. The $G_2$
alternative over a larger field was already excluded, and its
field-two exceptions have dimension six. The other exceptional
field-two cases have dimension four. Here $N\ge14$.
The classification therefore forces the image on $W$ to be
$\mathrm{SL}_N(2)=\mathrm{GL}_N(2)$.

For even $d$ this finishes the geometric assertion. For odd $d$
the kernel of the image in $\mathrm{GL}(W)$ is an invariant
subgroup of $\operatorname{Hom}(W,\langle\ell\rangle)$.
This is the irreducible dual natural module, so the kernel is
zero or the whole module. If it were zero, the group would be a
complement. Lemma~\ref{lem:char2-group-inputs} makes every such
complement conjugate to the block diagonal group. That group has
two separate orbits on $W_0\setminus\{0,\ell\}$, contradicting
the theta-orbit theorem. Thus the kernel is the whole module,
giving the full stabilizer and completing the induction.

The distinguished odd-degree covector is defined over the original
perfect field, as proved in the explicit Hasse--Witt construction.
Arithmetic monodromy is contained in its stabilizer and contains
the geometric group just determined; in even degree it is
contained in the full general linear group. Thus the two images
are equal in both cases. Irreducibility and the determinant
assertions follow immediately; $\mathbf F_2^\times$ is trivial.
\end{proof}

% Agent /root/monodromy; complete argument, 19 November 2026.
% Independently audited in full by /root, including the opened classification source.
\subsection{Plane curves in odd characteristic: the full group in every degree}
\label{sec:odd-plane-all-degrees}
\status{Proved; independently audited by /root, including the primary classification statement.}

Throughout this section $p$ is odd, $k$ is an algebraic closure of
$\mathbf F_p$, $d\geq4$, and $g=(d-1)(d-2)/2$.  Let $W_d$ be the
Hasse--Witt invertibility open in the projective space of degree-$d$
plane equations, including singular equations, and let $U_d\subset W_d$
be its smooth open.  The vector bundle $R^1\pi_*\mathcal O$ has constant
rank $g$ on the entire projective parameter space.  Its Lang fixed
sheaf $V$ is finite locally constant of rank $g$ on $W_d$.
The map $\pi_1(U_d)\to\pi_1(W_d)$ is surjective.  Thus computing its
image on $W_d$ computes the desired geometric image on $U_d$.
Choose as base point a union of $d$ lines with no triple intersection.
All groups below act on this same fibre, after choosing paths in the
specified connected parameter spaces.  This is permissible because
every boundary family used below is contained in $W_d$.

\subsubsection{The line-arrangement representation}

Put $P=\mathbf F_p^d$, with its permutation basis $e_1,\ldots,e_d$,
$u=e_1+\cdots+e_d$, and $W=P/\langle u\rangle$.  The graph of the
line arrangement is the complete graph $K_d$.  Its first cochain space
is $\bigwedge^2P$, with $e_i\wedge e_j$ representing an oriented edge.
The coboundaries are $u\wedge P$.  Consequently
\begin{equation}\label{eq:odd-arrangement-module}
 V=H^1(K_d,\mathbf F_p)
   =\bigwedge^2P/(u\wedge P)=\bigwedge^2W.
\end{equation}
This is cohomology, rather than the dual cycle space.  The coherent
normalization sequence gives the same identification after tensoring
with $k$, with bijective Frobenius on all the graph terms.

The space of ordered general-position lines is irreducible; its
quotient by the freely acting group $S_d$ is the unordered arrangement
space.  Its ordered cover is connected, so its permutation monodromy
is $S_d$.  Over the ordered cover the graph local system is constant.
Thus the image $G\subseteq\mathrm{GL}(V)$ contains exactly the
permutation action in \eqref{eq:odd-arrangement-module}.

Let $U=\{(x_i)\in P:\sum x_i=0\}$, let $D$ be its image in $W$, and
put $A=\bigwedge^2D\subseteq\bigwedge^2W$.

\begin{lemma}\label{lem:odd-wedge-submodules}
Every nonzero $S_d$-submodule of $\bigwedge^2W$ contains $A$.
If $p\nmid d$, then $A=V$ and $V$ is simple.  If $p\mid d$, then
$A$ is the unique proper nonzero submodule, and $V/A\simeq D$ is
simple of dimension $d-2$.
\end{lemma}

\begin{proof}
For four distinct indices $i,j,k,l$, put $r=e_i-e_j$ and $s=e_k-e_l$.
The group-algebra operator
\[
 (1-(ij))(1-(kl))
\]
on $\bigwedge^2W$ has image $\langle r\wedge s\rangle$ and is the
rank-one map obtained by evaluating the functional
$(x_i-x_j)\wedge(x_k-x_l)$ and multiplying by $r\wedge s$.
The latter vector is nonzero: the line $\langle u\rangle$ does not
meet $\langle r,s\rangle$, including when $d=4$, since $p\ne2$.

The vectors $r\wedge s$ for disjoint pairs span $\bigwedge^2U$ before
passage to $W$, and hence span $A$ afterwards.  Explicitly, on four
indices take $f_i=e_i-e_l$, $f_j=e_j-e_l$, $f_k=e_k-e_l$.  The three
perfect matchings give the columns
\[
 \begin{pmatrix}0&1&1\\1&0&-1\\-1&1&0\end{pmatrix}
\]
in the basis $f_i\wedge f_j,f_i\wedge f_k,f_j\wedge f_k$; their
determinant is $2$.  Every pair of difference-basis vectors belongs to
such a four-index set.  The same argument in the dual permutation
space shows that the corresponding functionals span
$\bigwedge^2W^*$, because $W^*$ is the sum-zero subspace of $P^*$.
For every nonzero vector of $V$, therefore, one of these rank-one
operators has nonzero value.  Its entire $S_d$-orbit spans $A$.
This proves the first assertion and also proves that $A$ is simple.

If $p\nmid d$, then $U\to W$ is an isomorphism.  If $p\mid d$,
$D=U/\langle u\rangle$ has dimension $d-2$, and the sum functional
$\sigma$ descends to $W$ with kernel $D$.  Contraction with $\sigma$
gives the exact sequence
\[
 0\longrightarrow\bigwedge^2D\longrightarrow\bigwedge^2W
 \xrightarrow{\ \iota_\sigma\ }D\longrightarrow0.
\]
The module $D$ is simple: a nonzero class has a representative whose
coordinates are not all equal; applying $1-(ij)$ for two unequal
coordinates produces a nonzero multiple of $e_i-e_j$, and the
translates of this difference span $D$.  Every proper nonzero
submodule of $V$ contains $A$ and has zero image in the simple
quotient unless it is all of $V$.  This proves uniqueness.
\end{proof}

\begin{proposition}\label{prop:odd-plane-irreducible}
The geometric monodromy group $G$ is irreducible on $V$.
\end{proposition}

\begin{proof}
Only the case $p\mid d$ needs another family.  Take a smooth ordinary
curve of degree $d-2$ and a smooth conic meeting it transversely in
$m=2d-4$ points.  The graph contribution is a sub-local-system
\[
 B=\mathbf F_p^m/\langle(1,\ldots,1)\rangle\subset V
\]
of rank $m-1=2d-5$.  Its node monodromy is the full symmetric group
$S_m$.  To verify this, fix the conic.  Restriction of degree-$(d-2)$
plane forms to it is surjective onto binary forms of degree $m$.
Over the squarefree binary forms, the ordered-root cover is connected.
Pulling it back gives an affine bundle after choosing a coefficient
chart (equivalently, an open of a projective bundle).  Imposing
smoothness and ordinarity of the plane curve removes an open subset
from this irreducible total space, and the open is nonempty.  The
ordered-root cover remains connected.  Thus the induced graph image
is $S_m$.

Here $p\nmid m$, since $p\mid d$ and $p$ is odd.  The deleted
permutation module $B$ is simple, by the same difference argument
used above.  If the only possible proper $G$-invariant subspace $A$
were invariant, the map $B\to V/A$ would be zero: its source is a
simple module of dimension $2d-5>d-2=\dim(V/A)$.  Hence $B\subset A$.

This inclusion can be tested at the common line arrangement.  Allow
the conic to become the union of the first two lines and the other
curve to become the remaining $d-2$ lines.  Keep the intersections
between the two subcurves transverse, disjoint from their internal
nodes.  Their ordinary coherent cohomology and the graph sequence
extend over an open containing both this point and the smooth
subcurve locus.  At this arrangement $B$ is the kernel of restriction
to the complete subgraph on vertices $3,\ldots,d$.  In particular the
edge cochain $e_1\wedge e_3$ represents an element of $B$.  But
\[
 \iota_\sigma(e_1\wedge e_3)=e_3-e_1\ne0\quad\hbox{in }D,
\]
so it does not belong to $A$.  This contradiction proves the claim.
No specialization through a drop of the Hasse--Witt rank has been
used.
\end{proof}

\subsubsection{Quartic boundaries and an elementary cocycle calculation}

\begin{lemma}\label{lem:sl3-natural-h1-odd}
For every odd prime $p$, the natural three-dimensional module $Q$
and its dual satisfy
$H^1(\mathrm{SL}_3(\mathbf F_p),Q)=
H^1(\mathrm{SL}_3(\mathbf F_p),Q^*)=0$.
\end{lemma}

\begin{proof}
Let $c$ be a cocycle.  The diagonal sign subgroup of determinant one
has order four, prime to $p$, so subtracting a coboundary makes $c$
zero on it.  Each $t_{ij}(a)=1+aE_{ij}$ commutes with the sign matrix
which is $-1$ in positions $i,j$ and $+1$ in the remaining position
$k$.  The cocycle relation then forces
$c(t_{ij}(a))\in\langle e_k\rangle$.
Write $x=t_{12}(a)$, $y=t_{23}(b)$, with $a,b\ne0$, and
$c(x)=Ae_3$, $c(y)=Be_1$.  Since $[x,y]=t_{13}(ab)$, the cocycle
identity gives
\[
 c([x,y])=-bA e_2-abA e_1.
\]
Its left side belongs to $\langle e_2\rangle$, hence $A=0$.
Permuting the three indices proves vanishing on every elementary
transvection.  They generate $\mathrm{SL}_3(\mathbf F_p)$, so the
normalized cocycle is zero.  Applying the inverse-transpose
automorphism proves the dual assertion.
\end{proof}

For any chosen four of the $d$ lines, fix the other $r=d-4$ lines and
vary an ordinary quartic, transverse to the fixed lines and avoiding
their pairwise intersections.  Include ordinary singular quartics
whenever the total degree-$d$ curve stays in $W_d$.  This connected
open contains the original arrangement.  Mayer--Vietoris for the
quartic and the fixed lines gives
\begin{equation}\label{eq:odd-quartic-boundary}
 0\longrightarrow T\longrightarrow V\longrightarrow Q
 \longrightarrow0,\qquad \dim Q=3,
\end{equation}
where $Q$ is the quartic Hasse--Witt local system.  Its image is
$\mathrm{GL}_3(\mathbf F_p)$ by the canonical quartic moduli argument
and Achter--Pries, Proposition~4.8, already checked above.  The image
on $T$ is solvable: its only moving graph edges are the four
intersections with each fixed line, so it factors through a subgroup
of $(S_4)^r$.  The nodes between the fixed lines are fixed.

Let $B$ be the image of this boundary family and let $D_T$ be its
image on $T$.  The image of $B$ in
$\mathrm{GL}(Q)\times D_T$ is a subdirect product.  Goursat's lemma,
or the elementary description of a subdirect product as equality in
a common quotient, shows that it contains
$\mathrm{SL}(Q)\times1$: every common quotient is solvable, whereas
$\mathrm{SL}_3(\mathbf F_p)$ is perfect.  Let $B_0$ be the inverse
image of $\mathrm{SL}(Q)\times1$.  It is normal in $B$, fixes $T$
pointwise, and maps onto $\mathrm{SL}(Q)$.
Its kernel is an additive submodule
\begin{equation}\label{eq:odd-quartic-kernel}
 K\subseteq\operatorname{Hom}(Q,T)=Q^*\otimes T.
\end{equation}
The natural dual is absolutely simple, with endomorphism ring
$\mathbf F_p$.  A submodule of this direct sum of copies of $Q^*$
therefore has the form $K=Q^*\otimes T'$ for a subspace $T'\subseteq T$.
The subspace $T'$ is $B$-invariant because $K$ is normal in $B$.
If $K\ne0$, it contains rank-one maps $w\otimes\phi$ and hence
transvections $1+w\otimes\phi$ of $V$.  If $K=0$, $B_0$ is a
complement to $\operatorname{Hom}(Q,T)$ in the indicated semidirect
product.  Lemma~\ref{lem:sl3-natural-h1-odd}, applied to each copy of
$Q^*$, makes that complement conjugate to the block-diagonal
$\mathrm{SL}(Q)$.  In this case too it contains transvections of $V$.
For $d=4$ the statement simply recovers the full quartic group.

\subsubsection{The normal subgroup generated by transvections}

\begin{lemma}\label{lem:transvection-clifford-blocks}
Let an irreducible finite group $G\subseteq\mathrm{GL}(V)$ contain a
transvection, and let $N\triangleleft G$ be the subgroup generated by
all transvections in $G$.  Then $V$ is a direct sum of pairwise
nonisomorphic simple $N$-modules, and $G$ permutes these summands
transitively.  Every transvection has its image line in exactly one
summand and fixes all the others pointwise.
\end{lemma}

\begin{proof}
The sum of the simple $N$-submodules, $\operatorname{soc}_N(V)$, is
nonzero and $G$-stable, since $N$ is normal.  Irreducibility of $G$
therefore makes it all of $V$.  The canonical isotypic summands are
permuted transitively by $G$.  Write any one of them as $S^{\oplus e}$.
For an element of $N$, the rank of its difference from the identity
on this summand is a multiple of $e$.  Some transvection has nonzero
action on some summand, and its total rank is one, so $e=1$ there;
transitivity gives $e=1$ for all summands.  Thus the canonical
isotypic summands are themselves simple.  Each transvection is
block diagonal because it belongs to $N$, and its rank-one
difference is nonzero on only one block.  This proves the assertion.
(The same argument with the finite division endomorphism ring of a
block shows it is absolutely simple, but that strengthening is not
needed below.)
\end{proof}

For a four-element subset $I\subset\{1,\ldots,d\}$, let
$U_I\subset P$ be its three-dimensional sum-zero permutation
subspace, supported on $I$.  Because $4\ne0$ in $\mathbf F_p$,
$U_I$ embeds in $W$.  Set
\[
 S_I=\bigwedge^2 U_I\subset V.
\]
The subgroup $S_I^{\rm perm}\simeq S_4$ permuting these four lines
lies in the corresponding quartic boundary group $B$.  Restricted
to it, the decomposition of the permutation space gives
\begin{equation}\label{eq:s4-quartic-decomposition}
 W\simeq U_I\oplus\mathbf F_p^{d-4},\qquad
 V\simeq S_I\oplus U_I^{\oplus(d-4)}
             \oplus\mathbf F_p^{\binom{d-4}{2}}.
\end{equation}
Both $U_I$ and $S_I$ are simple.  For the former use the elementary
deleted-permutation argument with $p\nmid4$; for the latter use
Lemma~\ref{lem:odd-wedge-submodules} with four indices.  They are
nonisomorphic: a transposition has trace $1$ on $U_I$ and trace $-1$
on $\bigwedge^2U_I$, distinct in every odd characteristic.  Neither
is trivial.  Thus \eqref{eq:s4-quartic-decomposition} is a semisimple
decomposition even when $p=3$, and $S_I$ occurs with multiplicity one.
The quotient $Q$ in \eqref{eq:odd-quartic-boundary} is isomorphic to
$S_I$ as an $S_4$-module, and $T$ contains no constituent of this type.

\begin{proposition}\label{prop:odd-transvection-normal-irreducible}
The normal subgroup $N\triangleleft G$ generated by all transvections
is irreducible on $V$.
\end{proposition}

\begin{proof}
Use the canonical $N$-block decomposition of
Lemma~\ref{lem:transvection-clifford-blocks}.  We first prove that for
each quartet $I$ its subspace $S_I$ lies entirely in one block.

If the kernel $K$ of its boundary group $B_0$ is zero, the cocycle
calculation makes $B_0$ a conjugate block-diagonal $\mathrm{SL}_3$.
Its commutator subspace $P_I=[B_0,V]$ is therefore three dimensional,
simple as a $B_0$-module, and maps isomorphically to $Q$.  Normality
of $B_0$ in $B$ makes $P_I$ stable under the quartet's $S_4$.
The unique-copy statement in \eqref{eq:s4-quartic-decomposition}
forces $P_I=S_I$.  This copy of $\mathrm{SL}_3$ is generated by
rank-one transvections of $V$, so $B_0\subseteq N$.  Consequently
each $N$-block is $B_0$-stable and
$[B_0,V]=\bigoplus_j[B_0,V_j]$.  Simplicity of $P_I$ puts it in one
$N$-block.

If $K\ne0$, write it as $Q^*\otimes T'$, with $T'\ne0$, and fix
one nonzero $\phi\in Q^*$.  For every nonzero $w\in T'$, the
transvection $1+w\otimes\phi$ is supported on a single $N$-block.
The functional here means the pullback of $\phi$ along $V\to Q$.
It is nonzero on that block and zero on every other block.  Since
the same nonzero functional is used for all $w$, that block must
be the same block $L$ for every $w$.  Thus $T'\subset L$.
Now fix one nonzero $w$ and vary $\phi$.  Every $\phi\in Q^*$
vanishes on every other $N$-block.  It follows that $L\to Q$ is
surjective.  Since $T'$ is stable under the quartet's $S_4$, this
$S_4$ must preserve $L$: its translate is another $N$-block sharing
the nonzero subspace $T'$.  As an $S_4$-submodule of the semisimple
module \eqref{eq:s4-quartic-decomposition}, $L$ is semisimple.
Its surjection onto $Q\simeq S_I$ forces it to contain the unique
copy $S_I$ in $V$.  This proves the claim in both cases.

When two quartets have three common indices, their subspaces $S_I$
intersect in the nonzero line $\bigwedge^2U_{I\cap J}$.
The graph of four-element subsets joined by such overlaps is
connected.  Hence all $S_I$ belong to the same $N$-block $L$.
Their span is $A=\bigwedge^2D$: the spanning disjoint-difference
wedges in Lemma~\ref{lem:odd-wedge-submodules} are supported on
quartets.  Thus $A\subset L$.

The subspace $A$ is $S_d$-stable.  Every translate of $L$ by $S_d$
is an $N$-block containing $A$, hence is $L$ itself.  Thus $S_d$
preserves $L$ and the sum $L'$ of the other blocks.  If $L'\ne0$,
Lemma~\ref{lem:odd-wedge-submodules} forces $A\subset L'$, which
contradicts $A\subset L$ and $L\cap L'=0$.  There are no other
blocks, and $N$ is irreducible.  The base case $d=4$ was already
settled by the quartic monodromy theorem.
\end{proof}

\subsubsection{Classification, exclusion of the alternating form, and determinant}

McLaughlin's transvection theorem, in the precise form of
Kantor--Seress (2001), Theorem~3.6, p.~22 \cite{KantorSeress2001},
states that for $n\geq3$
an irreducible subgroup of $\mathrm{SL}_n(\mathbf F_q)$ generated by
full transvection groups is $\mathrm{SL}_n(\mathbf F_q)$ or, when
$n$ is even, $\mathrm{Sp}_n(\mathbf F_q)$; the additional alternatives
in that theorem occur only when $q=2$.  Here $q=p$ is prime, and
the powers of $1+v\phi$ are precisely $1+a v\phi$, $a\in\mathbf F_p$.
Thus every generating transvection supplies its full transvection
group.  The theorem applies to $N$, which has determinant one and
is irreducible by Proposition~\ref{prop:odd-transvection-normal-irreducible}.
It leaves $N=\mathrm{SL}_g(\mathbf F_p)$ or $N=\mathrm{Sp}_g(\mathbf F_p)$.

The symplectic possibility is impossible.  If $N$ is symplectic,
normality puts $G$ in its group of similitudes: the invariant
alternating form is unique up to scalar.  For any quartet, the
multiplier on $B_0$ is trivial.  Indeed it is trivial on its
$p$-group kernel $K$ and on its perfect quotient
$\mathrm{SL}_3(\mathbf F_p)$.  Therefore $B_0$ preserves the
nondegenerate alternating form.  Since it fixes $T$ pointwise and
has fixed-point-free natural quotient $Q$, its fixed space is
$V^{B_0}=T$.  For a group preserving a nondegenerate bilinear form,
\[
 [B_0,V]=(V^{B_0})^\perp=T^\perp.
\]
The latter space has dimension three.  Its image in $Q$ is
$[B_0,Q]=Q$, so it maps isomorphically to $Q$ and meets $T$ trivially.
The alternating form is consequently nondegenerate on this
three-dimensional space $T^\perp$, a contradiction.  We have proved
\begin{equation}\label{eq:odd-plane-sl-contained}
 \mathrm{SL}_g(\mathbf F_p)\subseteq G.
\end{equation}

It remains to show that the geometric determinant is all of
$\mathbf F_p^\times$.  Take a Legendre elliptic cubic $E_t$, near a
supersingular parameter $t_0$, and add $d-3$ fixed lines chosen in
general position relative to $E_{t_0}$.  The Hasse polynomial of
the Legendre family has a simple root at $t_0\notin\{0,1\}$, as
proved in the elliptic calculation in
\S\ref{subsec:plane-residues}.  Over the strict local
trait at $t_0$, all graph contributions have bijective Frobenius
and determinant a unit.  The normalization sequence therefore
makes the total Hasse determinant a unit times the elliptic Hasse
invariant.  It has valuation exactly one at $t_0$.
On the punctured trait the total curve is ordinary, hence belongs
to $W_d$.  The determinant Lang torsor is
$z^{p-1}=h^{-1}$ by \eqref{eq:determinant-kummer}.  A valuation-one
$h$ gives tame inertia image of exact order $p-1$.  This is
geometric inertia, since the residue field is algebraically closed.
Thus $\det G=\mathbf F_p^\times$.  Combining this with
\eqref{eq:odd-plane-sl-contained} proves the theorem.

\begin{theorem}\label{thm:plane-odd-all-groups}
Let $p$ be odd and let $U_d$ be the ordinary smooth plane-curve locus
of degree $d$ over any perfect field $k$ of characteristic $p$.  At a
geometric generic point put
$V=H^1_{\mathrm{et}}(C_{\bar\eta},\mathbf F_p)$, of dimension
$g=(d-1)(d-2)/2$.  For every $d\geq3$, both the geometric and the
arithmetic monodromy images on $V$ are
$\mathrm{GL}_g(\mathbf F_p)$, with full determinant character.
For $d=1,2$ both images are trivial on the zero vector space.
For positive rank the representation is irreducible.
\end{theorem}

\begin{proof}
The preceding argument proves the geometric assertion for $d\geq4$
over $\overline{\mathbf F}_p$.  The connected finite etale frame cover
which it identifies is geometrically connected and remains connected
after every algebraically closed field extension.  Since the family
and its Frobenius local system are defined over $\mathbf F_p$, this
gives the same geometric image over the algebraic closure of any
perfect field of characteristic $p$.
The arithmetic image contains that image and is a subgroup of the
same finite general linear group, so it is equal to it.  The cubic
case follows from the same Legendre determinant calculation, since
its rank is one.  Degrees one and two have rank zero.
\end{proof}

% /root/monodromy, 19 November 2026.
% Independently audited in full by /root, including the determinant seeds.
\subsection{Hypersurfaces in projective space in odd characteristic}
\label{subsec:projective-space-odd-monodromy}
\status{Proved; independently audited by /root, including both determinant seeds.}

Fix an odd prime $p$ and an algebraically closed field $k$ of
characteristic $p$.  For degree $d$ hypersurfaces in $\mathbf P^n$,
let $W_{n,d}$ be the Hasse--Witt invertibility open in the full
projective space of equations, and let
\[
 h_{n,d}=\binom{d-1}{n},\qquad
 V_{n,d}=\left(H^{n-1}(\mathcal O_Y),F=1\right).
\]
The binomial coefficient is zero for $d\leq n$.
The full coherent bundle has this constant rank for all nonzero
equations, including singular equations.  A general-position union
of $d$ hyperplanes belongs to $W_{n,d}$, so this open is nonempty.
The desired smooth ordinary locus is a dense open in $W_{n,d}$;
passing to it does not change the finite monodromy image.
All the comparison families below therefore take place inside
one explicit constant-rank Lang local system, except for a
valuation-one Hasse degeneration used solely to compute its tame
determinant inertia.

\subsubsection{A connectedness lemma for linear projection}

\begin{lemma}\label{lem:linear-projection-monodromy}
Let $U$ be a nonempty open in $\mathbf P^N_k$, and let
$T\to U$ be a connected finite etale cover.  Let
$f:\mathbf P^N\dashrightarrow\mathbf P^M$ be a linear projection
whose center $Z$ meets $U$.  The restriction of $T$ to the
geometric generic fibre of $f|_{U\setminus Z}$ is connected.
Consequently a finite local system on $U$ retains its full
monodromy after restriction to that geometric generic fibre.
\end{lemma}

\begin{proof}
The open $U$ is smooth and integral, so $T$ is smooth and integral.
Let $B=\operatorname{Bl}_Z(\mathbf P^N)$ and put
$T'=T\times_U B_U$, where $B_U$ is the inverse image of $U$.
This is the blowup of $T$ in the smooth etale inverse image of
$Z\cap U$; in particular it is smooth and integral.  The resolved
projection gives a morphism $T'\to\mathbf P^M$.
Choose $z\in Z\cap U$.  The exceptional fibre above $z$ is
$E_z\simeq\mathbf P^M$, mapping isomorphically to the target.
Its pullback to $T'$ is the disjoint union of copies of $E_z$,
since it is the product of $E_z$ with the finite reduced set $T_z$.

Put $K=k(\mathbf P^M)$, and let $L$ be the algebraic closure of $K$
inside $k(T')$.  It is a finite extension of $K$.  Let $C$ be the
normalization of $\mathbf P^M$ in $L$; it is finite and integral.
There is a morphism $T'\to C$ over $\mathbf P^M$.  Indeed, on an
affine open of the target, each element of its integral closure
in $L$ is integral over the local ring at every point of the
normal scheme $T'$ lying above this open.  It therefore belongs
to each such local ring and defines a regular function there.
These morphisms glue.  One copy of $E_z$ in $T'$ consequently
gives a section $\mathbf P^M\to C$ of the finite map.
Its image is closed, has dimension $M$, and lies in the integral
scheme $C$ of dimension $M$; it is all of $C$.  Thus $C\simeq
\mathbf P^M$ and $L=K$.

On the nonempty open $T|_{U\setminus Z}$, the map to $\mathbf P^M$
is a composition of a finite etale map and a smooth linear
projection.  Hence $k(T')/K$ is separably generated.  The equality
$L=K$ implies that this extension is regular, so its geometric
generic fibre is integral.  Removing the exceptional divisor
does not alter that assertion.  For the local-system consequence,
apply the proved assertion to a connected component of its finite
frame cover, which is a torsor for its full monodromy image.
\end{proof}

\begin{corollary}\label{cor:restriction-independent-monodromy}
For $e\geq n+1$, fix a hyperplane $H\subset\mathbf P^n$.
The local system $V_{n,e}$ retains its full geometric monodromy
when the restriction of the equation to $H$ is prescribed to be
geometric generic.  On the parameter space where both the
equation and its restriction have invertible Hasse--Witt maps,
the joint monodromy image on $V_{n,e}$ and $V_{n-1,e}$ is the
product of their two full images.
\end{corollary}

\begin{proof}
Restriction of equations is a linear projection.  Its center
consists of equations divisible by the fixed equation of $H$.
It meets $W_{n,e}$: take $H$ and $e-1$ other general-position
hyperplanes.  Lemma~\ref{lem:linear-projection-monodromy} shows
that the joint image contains the first full group times $1$.
Its projection to the second group is surjective, since the
restriction map has geometrically integral generic fibre, an
open of an affine space.  The joint image is therefore the
product.  Removing further nonempty opens does not change
these finite images.
\end{proof}

\subsubsection{The arrangement module and independent quadratic reflections}

Take a general-position union of $d$ hyperplanes as base point.
Put $P=\mathbf F_p^d$, $u=\sum_i e_i$, $W=P/\langle u\rangle$,
$U=\ker(\sigma:P\to\mathbf F_p)$ with $\sigma(x)=\sum_i x_i$,
and $D=\operatorname{im}(U\to W)$.  The Cech resolution of the
union, whose strata are projective spaces, identifies its top
coherent cohomology and its Lang fibre with
\begin{equation}\label{eq:pn-arrangement-module}
 V_{n,d}=\bigwedge^nP/(u\wedge\bigwedge^{n-1}P)
       =\bigwedge^nW.
\end{equation}
The connected ordered-hyperplane cover of the unordered
arrangement locus gives the full permutation action of $S_d$
on this space.  Put $A=\bigwedge^nD$.

\begin{lemma}\label{lem:odd-general-hook-module}
Assume $1\leq n\leq d-2$.  If $p\nmid d$, then
$\bigwedge^nW$ is simple as an $S_d$-module.  If $p\mid d$, then
$A=\bigwedge^nD$ is its unique proper nonzero submodule; both
$A$ and the quotient $\bigwedge^{n-1}D$ are simple.
\end{lemma}

\begin{proof}
First suppose $2n\leq d$.  For $n$ pairwise disjoint transpositions
$\tau_j=(a_jb_j)$, the group-algebra operator
$\prod_j(1-\tau_j)$ has rank one on $\bigwedge^nW$.  Its image
is the wedge of the $n$ differences $e_{a_j}-e_{b_j}$, and its
functional is the wedge of the corresponding coordinate
differences.  This follows by simultaneously decomposing $W$
into the $n$ negative eigenlines and their common positive
eigenspace; $2$ is invertible.

These matching wedges span $\bigwedge^nU$.  To check this without
a modular representation-theory citation, use the difference
basis $f_i=e_i-e_o$ and the identity, for four distinct indices,
\[
 2f_a\wedge f_b
 =(e_a-e_b)\wedge(e_c-e_o)
  +(e_a-e_c)\wedge(e_b-e_o)
  +(e_a-e_o)\wedge(e_b-e_c).
\]
Beginning with a wedge of $n$ differences sharing $o$, repeatedly
introduce one unused index $c$.  Each application leaves one
fewer difference sharing $o$ and produces one disjoint pair.
There are $n-1$ available new indices because $d\geq2n$.
The resulting wedges are perfect matchings on $2n$ indices.
The same proof in the dual permutation space shows that the
matching functionals span $\bigwedge^nW^*$.
Thus every nonzero submodule contains a matching image and hence
all of $A$.  This proves simplicity of $A$, and proves simplicity
of the whole space when $p\nmid d$.

When $p\mid d$, the dot product gives a nondegenerate invariant
form on $D=U/\langle u\rangle$.  The same matching argument for
$n-1$ shows that $\bigwedge^{n-1}D$ is simple: matching images
span it, and their functionals span its dual under this form.
Contraction with the sum functional gives an exact sequence
\[
 0\longrightarrow\bigwedge^nD\longrightarrow\bigwedge^nW
 \longrightarrow\bigwedge^{n-1}D\longrightarrow0.
 \tag{*}
\]
Every nonzero submodule contains its kernel, and a submodule
which also surjects to the simple quotient is the whole space.
This proves the unique-submodule assertion in the stated range.

For $2n>d$, put $r=d-1-n$, so $1\leq r<d/2$.
The natural exterior pairing gives
$(\bigwedge^nW)^*\simeq\bigwedge^rW\otimes\operatorname{sgn}$.
If $p\nmid d$, the already proved simplicity for $r$ gives the
claim.  If $p\mid d$, duality preserves the property of having
exactly one proper nonzero submodule.  The displayed subspace
$\bigwedge^nD$ is nonzero and proper, so it must be that unique
submodule.  Its simplicity and that of the quotient follow
from the corresponding two simple factors for $r$.
\end{proof}

\begin{lemma}\label{lem:quadric-reflections}
For $n\geq2$ and $d\geq n+2$, the geometric group on
$V_{n,d}$ contains commuting rank-one reflections
$R_I$, indexed by $(n-1)$-subsets
$I\subset\{3,\ldots,d\}$.  At the arrangement base point,
their negative eigenlines are
\[
 \left\langle (e_1-e_2)\wedge\bigwedge_{i\in I}e_i\right\rangle,
\]
and their product is the permutation transposition $(12)$.
Here a rank-one reflection has eigenvalue $-1$ once and
eigenvalue $1$ on a codimension-one subspace.
\end{lemma}

\begin{proof}
Fix the last $d-2$ hyperplanes and replace the first two by a
variable quadric $Q$.  Require every positive-dimensional
quadric stratum to be connected with no higher coherent
cohomology, and every intersection with a line
$L_I=\bigcap_{i\in I}H_i$ to consist of two distinct points.
This is a nonempty open containing a smooth general quadric
and the reducible quadric $H_1H_2$.  A positive-dimensional
quadric hypersurface has $H^0(\mathcal O)=k$ and
$H^{>0}(\mathcal O)=0$, so the Cech complex computes all the
cohomology from its $H^0$ terms.  Only its top term has
nonconstant point sets: the two points on each $L_I$.
Swapping one pair acts on cohomology as a rank-one reflection,
since every coboundary has equal values on that pair.

All these swaps occur independently.  The restriction from
quadratic forms to $L_I\simeq\mathbf P^1$ is surjective.  The
discriminant of a binary quadratic is an irreducible polynomial
in its three coefficients in odd characteristic.  Its pullback
to the space of all quadratic forms is still irreducible, and
the discriminants for distinct lines are nonassociate.  To see
the last assertion, choose a point of $L_I\setminus L_J$ and
two hyperplanes through that point which meet $L_J$ in distinct
points.  Their product is tangent to $L_I$ and transverse to
$L_J$.  Their classes in the rational function field
modulo squares are therefore independent, by valuations at
these distinct prime divisors.  The ordered-pair cover has
Galois group $(\mathbf Z/2)^{\binom{d-2}{n-1}}$.
Removing the indicated open conditions leaves it connected.

The pair differences at $Q=H_1H_2$ are precisely the displayed
cochains, using the refinement of the cover with component $Q$
by its two hyperplanes.  Interchanging these two hyperplanes
interchanges every pair and fixes every other Cech term.
The transposition $(12)$ is thus the product of all the $R_I$.
Every fibre here has bijective Frobenius on the full top
coherent cohomology, because all terms are finite sets with
ordinary constant Frobenius.  The entire comparison is inside
$W_{n,d}$.
\end{proof}

\begin{proposition}\label{prop:pn-odd-irreducible}
For $n\geq2$ and $d\geq n+2$, the geometric monodromy group
$G_{n,d}$ is irreducible.
\end{proposition}

\begin{proof}
Only $p\mid d$ needs consideration, by
Lemma~\ref{lem:odd-general-hook-module}.  In that case any
proper invariant subspace must be $A=\bigwedge^nD$.
The negative eigenline of every $R_I$ lies outside $A$:
contraction with $\sigma$ sends its displayed generator to
\[
 -(e_1-e_2)\wedge\iota_\sigma
                  \left(\bigwedge_{i\in I}e_i\right)\ne0.
 \tag{**}
\]
Indeed the involved $n$ vectors remain linearly independent
modulo $u$: at least one index is outside $\{1,2\}\cup I$.
If $A$ were invariant, each $R_I$ would consequently fix $A$
pointwise.  Their product $(12)$ would do so too.  But on $D$
this transposition has a one-dimensional negative eigenspace
and a positive eigenspace of dimension $d-3$; on
$\bigwedge^nD$ its negative eigenspace has dimension
$\binom{d-3}{n-1}>0$.  This contradicts pointwise fixation.
\end{proof}


\subsubsection{The degree $n+2$ seed}

We use the following elementary observation about the Cech
computations, also in the determinant calculation below.  A family
consisting of coordinate hyperplanes and one further hypersurface
has the usual exact alternating resolution by its intersections
provided the latter equation restricts to a nonzero equation on
every coordinate stratum.  Locally its vanishing equation and any
of the coordinate equations form a regular sequence.  The
alternating union complex for a regular sequence is exact, by
induction using $(f)\cap(g)=(fg)$ and its versions modulo the
other equations.  Thus smoothness of the further hypersurface is
unnecessary.  Higher coherent cohomology of any hypersurface
stratum is given by its usual projective-space exact sequence.

\begin{proposition}\label{prop:pn-odd-seed-sl}
For $n\geq3$ and $d=n+2$, the group $G_{n,d}$ contains
$\mathrm{SL}_{n+1}(\mathbf F_p)$.
\end{proposition}

\begin{proof}
Let $N$ be the normal subgroup of $G_{n,d}$ generated by all its
transvections.  Replace any four of the arrangement hyperplanes,
indexed by $I$, by a varying quartic $Q$, and fix the other
$n-2$ hyperplanes.  The Cech spectral sequence has precisely
the following positive-cohomology terms contributing to total
degree $n-1$:
the quartic plane curve obtained by intersecting all the fixed
hyperplanes, of rank three, and the $n-2$ quartic surface strata,
each of rank one.  Higher-dimensional quartic strata have no
positive coherent cohomology.  The $H^0$ row has no positive
cohomology because the nerve is a full simplex.  There are no
nonzero incoming or outgoing differentials for the displayed
terms.  We consequently obtain an exact sequence of Frobenius
modules, and of Lang local systems on the invertibility open,
\[
 0\longrightarrow Q_I\longrightarrow V_{n,n+2}
      \longrightarrow T_I\longrightarrow0,
 \qquad \dim Q_I=3,\quad\dim T_I=n-2.
 \tag{1}
\]
Here $T_I$ is the direct sum of the rank-one surface local systems.
The open where all these terms are ordinary contains the original
arrangement.  Its image on $Q_I$ is $\mathrm{GL}_3(\mathbf F_p)$:
restriction of quartics to the fixed plane is a linear projection
with affine generic fibre, and the ordinary plane-quartic group
is already proved.  The inverse image of its connected frame
cover remains integral; imposing the other nonempty open
conditions preserves connectedness.  The image on $T_I$ is
abelian, being a subgroup of $(\mathbf F_p^\times)^{n-2}$.

Apply the quartic-boundary argument of the preceding plane
section to the dual of (1).  Goursat's lemma gives a normal
subgroup $B_0$ fixing $T_I^*$ pointwise and mapping onto
$\mathrm{SL}(Q_I^*)$.  Its kernel is
$Q_I\otimes T'$ for a subspace $T'\subseteq T_I^*$.
If the kernel is nonzero, it contains rank-one transvections.
If it is zero, Lemma~\ref{lem:sl3-natural-h1-odd} splits the
extension and again gives transvections.  Thus $N$ is nontrivial.
By Proposition~\ref{prop:pn-odd-irreducible} and
Lemma~\ref{lem:transvection-clifford-blocks}, $V_{n,n+2}^*$ is a
direct sum of canonical simple $N$-blocks permuted transitively
by $G_{n,n+2}$.

In the arrangement fibre there is the $S_d$-equivariant
identification
\[
 V_{n,n+2}^* \simeq W\otimes\operatorname{sgn},
 \qquad d=n+2.
 \tag{2}
\]
For a quartet $I$, let $U_I$ denote its three-dimensional
sum-zero subspace in $P$.  Since $4$ is invertible, it embeds
in $W$.  The restriction of (2) to its $S_4$ is the direct sum
of the unique simple module
$S_I=U_I\otimes\operatorname{sgn}$, of dimension three, and
$d-4$ copies of the sign character.  This is an actual
semisimple decomposition in every odd characteristic.
The quotient $Q_I^*$ in the dual of (1) is this unique
three-dimensional constituent.

The two cases in the proof of
Proposition~\ref{prop:odd-transvection-normal-irreducible} now
apply verbatim, with this specified decomposition: if the kernel
is zero, the commutator space of the split $\mathrm{SL}_3$ is
$S_I$ and belongs to one $N$-block; if it is nonzero, the
rank-one maps with a common quotient functional put $T'$ in one
block, and varying the functional makes that block surject to
$Q_I^*$.  The block is $S_4$-stable because $T'$ is stable, so
it contains the unique constituent $S_I$.

Quartets with three common indices give $S_I$'s with nonzero
intersection.  Their connected overlap graph puts all these
subspaces in one block $L$.  They span
$D\otimes\operatorname{sgn}$.  Every translate of $L$ by $S_d$
contains this same nonzero subspace and must equal $L$.
Thus $S_d$ also preserves the sum of the remaining blocks.
If that sum were nonzero, the degree-one case of
Lemma~\ref{lem:odd-general-hook-module}, with the sign twist,
would force it to contain $D\otimes\operatorname{sgn}$ too.
This contradicts the direct sum.  Therefore $N$ is irreducible.

Kantor--Seress, Theorem~3.6 \cite{KantorSeress2001}, applies:
over the prime field, powers of each transvection supply its
full transvection group.  It leaves
$N=\mathrm{SL}_{n+1}(\mathbf F_p)$ or a symplectic group.
The latter is excluded by
Lemma~\ref{lem:pn-reflection-excludes-symplectic} below.
\end{proof}

\begin{lemma}\label{lem:pn-reflection-excludes-symplectic}
If an irreducible normal subgroup $N\triangleleft G$ is the
natural symplectic group in dimension $h\geq3$ over an odd prime
field, then $G$ cannot contain a rank-one reflection.
\end{lemma}

\begin{proof}
The space of $N$-invariant alternating forms is one dimensional,
so normality puts $G$ in the symplectic similitude group.
Let $R$ be such a reflection, with positive eigenspace of
dimension $h-1$ and negative eigenline.  If its multiplier
were different from one, its positive eigenspace would be
totally isotropic, contradicting $h-1>h/2$.
Its multiplier is therefore one.  The negative eigenline is
then orthogonal to the positive eigenspace, since $p\ne2$,
and to itself because the form is alternating.  It is in the
radical of the form, a contradiction.
\end{proof}

\subsubsection{Induction in the degree}

\begin{proposition}\label{prop:pn-odd-degree-induction}
Suppose the groups for degree $d-1$ in $\mathbf P^n$ and
$\mathbf P^{n-1}$ contain their full special linear groups,
where $n\geq3$ and $d\geq n+3$.  Then
$G_{n,d}$ contains $\mathrm{SL}_{h_{n,d}}(\mathbf F_p)$.
If in addition the first of these two groups has full
determinant, then so does $G_{n,d}$.
\end{proposition}

\begin{proof}
Use the boundary $Y=A\cup H$, with $A$ a degree-$(d-1)$
hypersurface and $H$ a fixed hyperplane.  The usual exact
sequence of the union gives, on the simultaneous ordinary open,
\begin{equation}\label{eq:pn-two-piece-boundary}
 0\longrightarrow S\longrightarrow V_{n,d}
       \longrightarrow Q\longrightarrow0,
 \qquad
 S=V_{n-1,d-1},\quad Q=V_{n,d-1}.
\end{equation}
There are no adjacent positive coherent groups, since a
hypersurface in projective space has no intermediate coherent
cohomology.  Put
\[
 r=\dim Q=\binom{d-2}{n},\qquad
 s=\dim S=\binom{d-2}{n-1}.
\]
Both are at least three in the present range.
Corollary~\ref{cor:restriction-independent-monodromy} shows
that the graded image is the product of the full groups.
Take the inverse image $B_0$ of
$L=\mathrm{SL}(S)\times\mathrm{SL}(Q)$.
Its kernel $K$ is an additive $L$-submodule of
$M=\operatorname{Hom}(Q,S)$.  This outer tensor product of
the two absolutely simple natural modules is absolutely simple,
so $K=0$ or $K=M$.

If $K=M$, its rank-one elements are transvections of $V$.
Let $N$ be their full normal closure in $G_{n,d}$ and use
the canonical $N$-blocks, justified by irreducibility of $G$.
For every $w\in S$ and $\phi\in Q^*$, the operator
$1+w\otimes\phi$ is a transvection when both are nonzero.
Fixing $\phi$ puts all of $S$ in one block; fixing a nonzero
$w$ and varying $\phi$ makes that block surject to $Q$.
It contains $S$ and surjects to $V/S$, so it is all of $V$.
Thus $N$ is irreducible.

If $K=0$, the complement is conjugate to the split block
diagonal group $L$.  Here is the needed cohomology computation,
which requires no vanishing theorem for either natural module:
\[
 H^1(\mathrm{SL}(S)\times\mathrm{SL}(Q),S\otimes Q^*)=0.
\]
The first factor has no invariants in $S$.  In the low-degree
Hochschild--Serre sequence the remaining term is
$H^1(\mathrm{SL}(S),S)\otimes Q^*$, whose invariants under the
second factor are zero.  This proves the displayed vanishing
and hence the conjugacy of complements.
The split $L$ is generated by rank-one transvections on $V$,
so is contained in $N$.  Every canonical $N$-block is therefore
an $L$-submodule of the semisimple, multiplicity-free module
$S\oplus Q$.  If there is more than one block, they are exactly
these two summands.  Transitivity of $G$ on the blocks forces
$r=s$, equivalently $d=2n+1$.

Even this last possibility is impossible.  A rank-one reflection
cannot interchange two subspaces of common dimension $r>1$:
its difference from the identity would have rank at least $r$
on either one.  Thus every reflection of
Lemma~\ref{lem:quadric-reflections} preserves both blocks.
Their product $(12)$ does so too.  Conjugating within the
arrangement $S_d$ shows that every transposition, hence all
of $S_d$, preserves both blocks.  This contradicts
Lemma~\ref{lem:odd-general-hook-module}: a simple module, or a
module with a unique proper nonzero submodule, has no such
nontrivial direct sum decomposition.  We have proved that
$N$ is irreducible also when $K=0$.

The transvection classification again gives $N=\mathrm{SL}(V)$
or a symplectic group.  The latter is excluded by
Lemma~\ref{lem:pn-reflection-excludes-symplectic}, using the
quadric reflection already present in $G$.
Finally the full product assertion for the graded image
contains the old group on $Q$ times $1$ on $S$.  Its lifts
have the same determinant as their action on $Q$, because
the extension kernel is unipotent.  Full determinant for
that old group therefore implies full determinant for $G$.
\end{proof}

\subsubsection{Determinants in the two initial degrees}

\begin{lemma}\label{lem:elliptic-fixed-three-points}
There is a family of plane cubics over a strictly henselian
smooth trait over $k$ whose special fibre is smooth supersingular,
whose Hasse invariant has valuation one, and whose restriction
to a fixed line is the constant squarefree binary cubic
$g=x_0x_1(x_0-x_1)$.  There are also smooth ordinary plane
cubics with this same restriction.
\end{lemma}

\begin{proof}
Take the Legendre family at a simple supersingular zero of its
Hasse polynomial, as proved in the earlier elliptic calculation.
Choose a line transverse to the special fibre and identify it
with $\mathbf P^1$.  Transversality persists on the trait, and
the three intersection points form a finite etale scheme.
They split after strict henselization.  The unique projective
linear transformation on the line taking their chosen order
to $0,1,\infty$ varies regularly with the trait parameter.
Lift this transformation to the ambient projective plane and
rescale the cubic equation by a unit.  Its restriction to the
line is then exactly $g$.  Isomorphisms and unit changes of
cohomology trivialization multiply the Hasse invariant by a
unit, so its valuation is still one.  The same construction
with any ordinary elliptic plane cubic gives the last assertion.
\end{proof}

\begin{proposition}\label{prop:pn-seed-full-determinant}
For every $n\geq2$, the geometric determinant images in degrees
$n+1$ and $n+2$ are the full group $\mathbf F_p^\times$.
\end{proposition}

\begin{proof}
For degree $n+1$, fix $n-2$ coordinate hyperplanes and let their
intersection be a plane.  Lift the cubic family of
Lemma~\ref{lem:elliptic-fixed-three-points} from this plane to
$\mathbf P^n$, for instance by its same equation independent
of the other coordinates, and take its union with the fixed
hyperplanes.  Their intersections are proper.  All
positive-dimensional cubic strata of ambient dimension at
least three have no positive coherent cohomology; the only
nonzero term in total degree $n-1$ is the elliptic curve term.
The $H^0$ nerve is a full simplex and has no higher cohomology.
Thus the rank-one total Hasse invariant is a unit times the
elliptic one and has valuation exactly one.
For $n=2$ this is simply the elliptic family itself.

For degree $n+2$, use coordinates
$x_0,x_1,y_1,\ldots,y_{n-1}$ and hyperplanes $H_i=(y_i=0)$.
Their common line is $\ell=\mathbf P(x_0,x_1)$.
Let $f_1(x_0,x_1,y_1)$ be the family in the lemma, and for
$i>1$ let $f_i(x_0,x_1,y_i)$ be a fixed ordinary plane cubic,
all restricting to $g$ on $\ell$.  Define explicitly
\[
 q=g+\sum_{i=1}^{n-1}(f_i-g),\qquad
 Y=(q\,y_1\cdots y_{n-1}=0)\subset\mathbf P^n.
 \tag{3}
\]
Restriction to the plane where all $y_j$, $j\ne i$, vanish
is exactly $f_i$.  Every coordinate stratum has a nonzero
restricted cubic equation, because its restriction to $\ell$
is $g$.  The alternating union resolution described above
is therefore exact, even if the higher-dimensional cubic
strata happen to be singular.

Those higher-dimensional cubic strata have no positive
coherent cohomology, and every positive-dimensional stratum
has $H^0(\mathcal O)=k$.  The common zero-dimensional stratum
consists of the three constant, distinct roots of $g$.
The $H^0$ row is the full-simplex complex except that its last
copy of $k$ is replaced by $k^3$; it contributes an ordinary
rank-two piece in degree $n-1$.  In the positive-cohomology
rows the $n-1$ elliptic curves contribute their rank-one
cohomology, again in total degree $n-1$.  There are no
possible nonzero differentials.  This gives a Frobenius-stable
filtration on the total cohomology with ordinary graph piece
of rank two and the $n-1$ elliptic pieces.  Its total rank is
$2+(n-1)=n+1=\binom{n+1}{n}$.
The Hasse determinant is a unit times the product of their
Hasse invariants.  Exactly $f_1$ has valuation one and all the
other factors are units, so the total determinant has
valuation one.

In both constructions the punctured trait lies in $W_{n,d}$.
The determinant Lang torsor is
$z^{p-1}=h^{-1}$, as in \eqref{eq:determinant-kummer}.
A valuation-one $h$ gives tame geometric inertia of order
$p-1$.  This proves the assertion.  The special fibre is
allowed to have smaller Hasse--Witt rank: only its valuation
is used, and no stable local system is extended across it.
\end{proof}

\begin{theorem}\label{thm:pn-odd-full-groups}
Let $p$ be odd, let $k$ be any perfect field of characteristic
$p$, and let $n\geq2$.  On the ordinary smooth degree-$d$
hypersurface locus in $\mathbf P^n_k$, the geometric and
arithmetic monodromy images on
$H^{n-1}_{\mathrm{et}}(Y_{\bar\eta},\mathbf F_p)$ are both
\[
 \mathrm{GL}_{\binom{d-1}{n}}(\mathbf F_p).
\]
For $d\leq n$ the vector space is zero and the image is trivial.
For $d=n+1$ it has rank one and full determinant.
For every $d\geq n+2$ it is irreducible with full determinant.
\end{theorem}

\begin{proof}
Dimension two is the already proved plane theorem.
Induct on $n\geq3$.  The zero-rank degrees need no argument,
and Proposition~\ref{prop:pn-seed-full-determinant} proves the
rank-one degree $n+1$.  In degree $n+2$,
Proposition~\ref{prop:pn-odd-seed-sl} and the determinant
proposition give the full general linear group.
Now induct on $d\geq n+3$ using
Proposition~\ref{prop:pn-odd-degree-induction}.
The two smaller groups are known by the dimension and degree
inductions, and the old determinant propagates as proved there.
This proves the geometric image over $\overline{\mathbf F}_p$.
The connected finite frame cover remains geometrically
connected under algebraically closed extension; since the
family is defined over $\mathbf F_p$, the result holds over
the algebraic closure of any perfect $k$.  The arithmetic
image contains the full geometric general linear group, so
is equal to it.
\end{proof}

\begin{corollary}\label{cor:pn-odd-all-residues-sharp}
Let $k$ be finite of odd characteristic $p$ and $n\geq2$.
For every $d\geq n+2$, every nonempty open in the smooth
degree-$d$ hypersurface locus of $\mathbf P^n_k$ contains,
over suitable finite extensions of $k$, hypersurfaces with
each of the $p$ point-count residues.
For $d\leq n$ every smooth hypersurface has residue $1$.
For $d=n+1$ the ordinary open misses residue $1$.
Thus $d=n+2$ is the exact threshold for the assertion involving
every nonempty open.
\end{corollary}

\begin{proof}
The finite-field coherent trace formula gives
\[
 \#Y(k')\equiv 1+(-1)^{n-1}
       \operatorname{Tr}(F_{k'}\mid H^{n-1}(Y,\mathcal O_Y)_{\rm ss})
       \pmod p,
\]
because a projective hypersurface has no other positive
coherent cohomology.  In the zero-rank range the trace is zero.
In rank one on the ordinary open it is nonzero.
In rank $h\geq2$, traces of invertible matrices take every
value in $\mathbf F_p$: the block
$\begin{pmatrix}0&-1\\1&a\end{pmatrix}$ has determinant one
and trace $a$, and adjoining an identity block allows any
specified total trace.  By the theorem the finite geometric
and arithmetic image is the full group.  The finite
Chebotarev theorem already cited in the plane calculation
therefore realizes each such conjugacy class in every
nonempty open after suitable finite extension of $k$.
The intersection of that open with the ordinary locus is
nonempty because the latter is dense.
\end{proof}

% /root/monodromy, 19 November 2026.
% Elementary all-dimension characteristic-two computation; root independently audited.
\subsection{Hypersurfaces in characteristic two: an elementary exact group calculation}
\label{subsec:pn-char2-exact}
\status{Proved; independently audited.}

Fix an algebraically closed field $k$ of characteristic two,
$n\geq2$, and $d\geq n+2$.  Put $m=d-1$ and
$h=\binom{m}{n}$.  Work on the full Hasse--Witt invertibility
open $W_{n,d}$ as in the preceding section.  Its Lang fibre
at a general-position hyperplane arrangement is
$V=\bigwedge^n W$, where
$W=\mathbf F_2^d/\langle e_0+\cdots+e_m\rangle$.
Use $e_1,\ldots,e_m$ as a basis of $W$, so
$e_0=e_1+\cdots+e_m$.
Let $b_A=\bigwedge_{i\in A}e_i$, indexed by the $n$-subsets
$A\subset\{1,\ldots,m\}$, be the exterior basis of $V$.
Signs disappear in this characteristic.

\begin{lemma}\label{lem:char2-independent-quadric-swaps}
The geometric monodromy contains, for each
$(n-1)$-subset $I\subset\{2,\ldots,m\}$, the transvection
\begin{equation}\label{eq:char2-quadric-transvection}
 R_I=1+w_I\phi_I,\qquad
 w_I=\sum_{\substack{a\in\{2,\ldots,m\}\\a\notin I}}b_{I\cup\{a\}},
 \quad
 \phi_I=b_{\{1\}\cup I}^*.
\end{equation}
It fixes every exterior basis vector except
$b_{\{1\}\cup I}$, which it sends to
$\sum_{a\notin I}b_{I\cup\{a\}}$.
The swaps for different $I$ are independent monodromy elements.
\end{lemma}

\begin{proof}
Replace $H_0H_1$ by a varying quadric and keep the remaining
hyperplanes fixed.  The same alternating resolution as in
Lemma~\ref{lem:quadric-reflections} has constant lower terms
and a pair of points for every line
$L_I=\bigcap_{i\in I}H_i$.  On an open containing the original
arrangement these pairs are distinct, all positive-dimensional
quadric strata are connected and have no higher coherent
cohomology, and the entire top coherent Frobenius is bijective.

We check independence of the point-pair covers in characteristic
two, working on the affine space of quadric equations.  This
cover is the pullback of the projective parameter cover, so
its full image is already a sufficient geometric lower bound.
On a chosen basis for a line $L$, the restricted quadratic
form is $a_LX^2+b_LXY+c_LY^2$.  Its ordered roots give the
Artin--Schreier extension
\[
 z^2+z=a_Lc_L/b_L^2
\]
of the rational function field of the space of quadrics.
The restriction map to any line is surjective, so
$a_L,b_L,c_L$ are independent linear functions.
Moreover $b_L$ is evaluation of the polar alternating form
on the Pluecker vector of $L$.  For distinct lines these
linear functions are nonassociate, by injectivity of the
Pluecker embedding.

At the prime divisor $b_{L_I}=0$, the displayed function for
$I$ has pole order two and leading residue $a_{L_I}c_{L_I}$.
This residue is not a square in the divisor's rational
function field, since $a_{L_I},c_{L_I}$ remain independent
variables there.  The corresponding functions for all
$J\ne I$ have no pole at this divisor.
No Artin--Schreier coboundary $q^2+q$ can have that leading
term: a pole of order two would require a pole of order one
for $q$ and would have square leading residue.  Thus no
nonempty sum of these functions is an Artin--Schreier
coboundary.  Their classes are independent, and the compositum
has Galois group $(\mathbf Z/2)^{\binom{d-2}{n-1}}$.
Passing to the nonempty ordinary and transverse open preserves
connectedness of this cover and hence all the independent swaps.

At the arrangement, swapping the two points on $L_I$ exchanges
the cochains $e_0\wedge e_I$ and $e_1\wedge e_I$, and fixes
all other top cochains not involving both $0$ and $1$.
These latter cochains span the arrangement cohomology; those
involving both are eliminated by the coboundary relation
$u\wedge\bigwedge^{n-1}P=0$.
In the chosen exterior basis the only changed column is
$b_{\{1\}\cup I}$, and its difference is
$(e_0+e_1)\wedge e_I=w_I$.
This proves \eqref{eq:char2-quadric-transvection}.
The vector $w_I$ is nonzero because
$d-n-1\geq1$, and its $\phi_I$ value is zero.
Thus the operator is a transvection.
\end{proof}

\begin{lemma}\label{lem:char2-coordinate-generation}
Let $G_0$ be generated by the arrangement permutation group
and the swaps of Lemma~\ref{lem:char2-independent-quadric-swaps}.
Writing $\ell=\sum_A b_A^*$, one has
\[
 G_0=
 \begin{cases}
 \mathrm{GL}_{h}(\mathbf F_2),&d-n-1\ \text{odd},\\
 \operatorname{Stab}_{\mathrm{GL}_h(\mathbf F_2)}(\ell),
       &d-n-1\ \text{even}.
 \end{cases}
\]
In the second case the stabilizer is
$\mathbf F_2^{h-1}\rtimes\mathrm{GL}_{h-1}(\mathbf F_2)$.
\end{lemma}

\begin{proof}
Permuting the labels $1,\ldots,m$ and fixing label $0$ supplies
the following operators for every $(n-1)$-subset $I$ and
every $j\notin I$:
\[
 T^I_j(b_{I\cup\{j\}})=S_I,\qquad
 S_I=\sum_{a\notin I}b_{I\cup\{a\}},
\]
with every other exterior basis vector fixed.  This operator
exchanges the two vectors $b_{I\cup\{j\}}$ and $S_I$ and
fixes each of the other vectors in this sum.
For $j,t\notin I$, $j\ne t$, the conjugate
$T^I_jT^I_tT^I_j$ is therefore the pure coordinate
transposition exchanging $b_{I\cup\{j\}}$ and
$b_{I\cup\{t\}}$ and fixing all other exterior basis vectors.
The graph of $n$-subsets joined when they share $n-1$ elements
is connected.  Its edge transpositions generate the full
coordinate symmetric group $S_h$, so $G_0$ contains that group.

Put $l=d-n-1$.  Each $T^I_j$ is a transvection whose image
vector has exactly $l$ nonzero exterior coordinates, all
different from its input coordinate.  Conjugating by $S_h$
gives every transvection
$1+v b_a^*$ with $v$ supported on any $l$ coordinates other
than $a$.  If $l=1$, these are all elementary transvections
and generate $\mathrm{GL}_h(\mathbf F_2)$.
If $l\geq2$, choose two such supports with $l-1$ common
coordinates and differing in any desired pair $b,c$, both
different from $a$.  This is possible because
$h=\binom{d-1}{n}\geq d-1\geq d-n+1=l+2$
for $n\geq2,d\geq n+2$.
The product, with the same input functional, is
\[
 1+(b_b+b_c)b_a^*.
 \tag{4}
\]
If $l$ is odd, adding pairs from an $l$-element support reduces
it to a singleton.  Thus once again all elementary
transvections occur and $G_0=\mathrm{GL}_h(\mathbf F_2)$.

Suppose $l$ is even.  Every generator fixes $\ell$:
the coordinate transvections change a column by an even-weight
vector, and an arrangement permutation fixes $\ell$ as well.
For the latter assertion, permutations of $1,\ldots,m$ act
by coordinate permutations, while $(01)$ is the product of
the commuting swaps in
Lemma~\ref{lem:char2-independent-quadric-swaps}.
These permutations generate the full arrangement $S_d$.
Consequently $G_0\subseteq\operatorname{Stab}(\ell)$.

Relabel the exterior basis as $b_1,\ldots,b_h$ and put
$f_i=b_i+b_h$ for $1\leq i<h$.
The $f_i$ form a basis of $\ker\ell$, with complement
$\langle b_h\rangle$.  Operators (4), with image $f_i$ and
input functional $b_j^*$ for $i\ne j<h$, give the elementary
matrices on $\ker\ell$ and fix $b_h$.  They generate its
full $\mathrm{GL}_{h-1}(\mathbf F_2)$.
The coordinate transposition $(i\,h)$ is
$1+f_i(b_i^*+b_h^*)$.  Multiplying it by the commuting
operators $1+f_i b_j^*$ for all $j\ne i,h$ gives
$1+f_i\ell$.  These generate all the translations of the
complement by $\ker\ell$.  Hence the full stabilizer is
contained in $G_0$, proving equality.
\end{proof}

\begin{lemma}\label{lem:pn-char2-canonical-covector}
If $d\equiv n+1\pmod2$ and $d\geq n+1$, the universal ordinary
Hasse--Witt local system has a nowhere-zero invariant
$\mathbf F_2$-valued covector, defined over $\mathbf F_2$.
\end{lemma}

\begin{proof}
Write a degree-$d$ equation as $F=\sum_{|\alpha|=d}c_\alpha x^\alpha$.
Use the standard top-cohomology basis indexed by positive
$(n+1)$-tuples $u$ with $|u|=d$.  The Hasse--Witt matrix is
$H_{uv}=c_{2v-u}$, with an out-of-range coefficient understood
as zero.  Put
\[
 w_u=\begin{cases}c_u,&\text{every coordinate of }u\text{ is odd},\\
                   0,&\text{otherwise}.\end{cases}
\]
Then
\[
 H^{\mathsf T}w=w^{[2]}.
 \tag{5}
\]
Indeed, in the $v$ coordinate the sum is
$\sum_{u\ {\rm odd}}c_{2v-u}c_u$.  Whenever the coefficient
is nonzero, $2v-u$ is also a positive all-odd tuple.
The involution $u\leftrightarrow2v-u$ cancels all paired
terms.  Its only possible fixed term is $u=v$, giving
$c_v^2$ when $v$ is all odd and zero otherwise.

There exists a positive all-odd tuple of sum $d$ under the
assumed parity and bound.  If every $w_u$ vanished, the row
of $H$ indexed by any such tuple would vanish: every
possibly nonzero coefficient in that row is itself indexed
by an all-odd tuple.  Thus $w\ne0$ wherever $H$ is invertible.
The map $x\mapsto w^{\mathsf T}x$ is Frobenius equivariant
by (5), and restricts on Lang fixed vectors to an
$\mathbf F_2$-valued covector.  It is nonzero because those
fixed vectors span the coherent fibre.
This construction is global on projective parameter space:
the coherent bundle is $\mathcal O(-1)$ tensored with the
fixed top-cohomology vector space, and the entries $w_u$
are degree-one functions of the equation.  They define an
honest morphism of this bundle to $\mathcal O$, not merely
a covector up to scalar.  Everything is defined over
$\mathbf F_2$.
\end{proof}

\begin{theorem}\label{thm:pn-char2-all-groups}
Let $k$ be any perfect field of characteristic two and
$n\geq2$.  Put $h=\binom{d-1}{n}$.  The geometric and
arithmetic monodromy groups on the ordinary smooth
degree-$d$ hypersurface locus of $\mathbf P^n_k$ are:
\[
 \begin{array}{c|c}
 \text{degree range}&\text{both monodromy images}\\ \hline
 d\leq n&1\text{ on the zero vector space}\\
 d=n+1&\mathrm{GL}_1(\mathbf F_2)=1\\
 d\geq n+2,\ d\equiv n\pmod2&\mathrm{GL}_h(\mathbf F_2)\\
 d\geq n+3,\ d\equiv n+1\pmod2&
    \operatorname{Stab}(\lambda)
       \simeq\mathbf F_2^{h-1}\rtimes
                         \mathrm{GL}_{h-1}(\mathbf F_2),
 \end{array}
\]
where $\lambda$ is the nonzero covector in
Lemma~\ref{lem:pn-char2-canonical-covector}.
In the last case $\ker\lambda$ is a proper nonzero invariant
hyperplane, so the representation is reducible.
\end{theorem}

\begin{proof}
The two initial ranks are zero and one as stated.
In all higher degrees
Lemma~\ref{lem:char2-coordinate-generation} gives the required
geometric lower bound.  In the parity with a canonical
covector, Lemma~\ref{lem:pn-char2-canonical-covector} gives
the matching upper bound, even for the arithmetic image.
For precision in comparing covectors: the subgroup $G_0$
has exactly the one nonzero invariant covector $\ell$.
The canonical $\lambda$ is nonzero and $G_0$-invariant,
so its value on the arrangement fibre is $\ell$.
Thus the two stabilizers are the same.
In the other parity the lower bound is the entire general
linear group.  The result over arbitrary perfect $k$
follows by geometric connectedness of the corresponding
finite cover under algebraically closed extension, exactly
as in the odd-characteristic theorem.
\end{proof}

\begin{corollary}\label{cor:pn-all-p-sharp-residue-threshold}
For every finite field $k$, every $n\geq2$, and every
$d\geq n+2$, each nonempty open in the smooth degree-$d$
hypersurface parameter space of $\mathbf P^n_k$ contains
hypersurfaces, over suitable finite extensions of $k$,
with every point-count residue modulo the characteristic.
The threshold $n+2$ is exact: for $d\leq n$ the only
residue is $1$, and for $d=n+1$ the ordinary open misses $1$.
\end{corollary}

\begin{proof}
Odd characteristic was settled in
Corollary~\ref{cor:pn-odd-all-residues-sharp}.  In characteristic
two, the full general linear group in rank at least two
has both trace values.  In the covector-stabilizer case,
the trace is $1+\operatorname{Tr}(B)$ with
$B\in\mathrm{GL}_{h-1}(\mathbf F_2)$; here $h-1\geq2$,
so again both values occur.  The same finite Chebotarev
and coherent trace argument gives both point-count
residues in every open.  The zero- and one-rank failures
are as in the preceding corollary.
\end{proof}

\par\smallskip\noindent\textbf{Finite checks.}
The retained program
\path{code/projective_quadric_generators.py} constructs
these exact exterior matrices over $\mathbf F_2$.
It checks that every swap has rank-one difference and
order two, that their product is the hyperplane
transposition, and that the common fixed covector space
has dimension one exactly in the indicated parity.
For the tested ranks at most fifteen it additionally
enumerates all nonzero vector orbits, obtaining either
$2^h-1$ or $2^{h-1}-1$ and $2^{h-1}$ as predicted.
These finite checks support the formulas; the uniform
group-generation proof above is independent of sampling.
